\chapter{Compléments aux chapitres du corps}
\label{chap:complements}

\section{Le marché français de la cyberassurance : trajectoire et segments}
\label{sec:ann-lucy}

% SOURCES-SCRIPTS: 63

% sec:lucy vivait dans 02_introduction.tex, devenu ORPHELIN le 19 septembre 2026 quand le
% chapitre 1 (02b_cadre_cyber_dora) a absorbe l'introduction. La lecture de marche LUCY y a
% ete reprise sous sec:marche-cyber-francais. Renvoi retabli le 21 septembre.
Complément du \S\ref{sec:marche-cyber-francais}.

\paragraph{Un marché qui s'élargit en baissant ses prix.} Entre $2024$ et $2025$, le nombre
d'entreprises assurées passe de $14\,124$ à $20\,996$, soit $+49\,\%$, quand les primes
souscrites reculent de $316{,}8$ à $305{,}9$~M\euro, soit $-3\,\%$ : l'élargissement du marché
s'est fait par la prime unitaire, qui perd $35\,\%$, et la franchise moyenne recule de
$149{,}9$ à $89$~k\euro. Dans le même temps, le nombre de sinistres indemnisés passe de $448$ à
$1\,251$.

\begin{figure}[htbp]
\centering
\figover{figures_externes/lucy_primes_sinistres_ratio.png}
\caption[Primes, sinistres et ratio du marché français de la cyberassurance]{Primes souscrites,
sinistres indemnisés et ratio sinistres sur primes du marché français de la cyberassurance,
$2019$ à $2025$, en millions d'euros. Reproduit du rapport LUCY 2026 de Nexialog Consulting
\citep{RapiorKaddouri2026}. L'étiquette du montant $2025$ des sinistres est recouverte par le
marqueur du ratio dans la source ; sa valeur est $83{,}2$~M\euro.}
\label{fig:lucy-sp}
\end{figure}

\noindent Le ratio sinistres sur primes remonte depuis trois exercices, de $12\,\%$ en $2023$ à
$17\,\%$ en $2024$ puis $27\,\%$ en $2025$, et l'écart double d'un exercice à l'autre, $5$ puis
$10$ points (figure~\ref{fig:lucy-sp}). Le marché a pourtant déjà connu un régime bien plus
dégradé, à $85\,\%$ en $2019$ et $169\,\%$ en $2020$ : le niveau de $2025$ est un point haut de
trois ans, non un point haut d'historique.

\noindent La dégradation n'est pas uniforme entre segments (figure~\ref{fig:lucy-divergence}). Le
ratio des entreprises de taille intermédiaire passe de $13\,\%$ en $2024$ à $42\,\%$ en $2025$,
quand l'indice de leur taux de prime recule de son sommet de $257$ en $2023$ à $178$ en $2025$ :
le prix baisse sur le segment dont le risque se dégrade. Ce segment a déjà connu un ratio de
$261\,\%$ en $2021$.

\begin{figure}[htbp]
\centering
\figover{figures_externes/lucy_divergence_prix_risque.png}
\caption[Divergence entre prix et risque par segment]{Ratios sinistres sur primes par segment
(barres, axe de gauche) et indice du taux de prime des entreprises de taille intermédiaire
(courbe, axe de droite, base $100$ en $2020$). Pour chaque exercice, les trois barres sont, de
gauche à droite, les grandes entreprises, les entreprises de taille intermédiaire et les
entreprises moyennes. Reproduit du rapport LUCY 2026 de Nexialog Consulting
\citep{RapiorKaddouri2026}.}
\label{fig:lucy-divergence}
\end{figure}

\vspace{2pt}
\noindent Un point de vocabulaire accompagne le tableau des deux régimes du \S\ref{subsec:lucy-freq-sev}. La \emph{fréquence} désigne un nombre de
sinistres rapporté au nombre d'assurés, tandis que le \emph{nombre} de sinistres est un total. Ce
sont deux grandeurs distinctes : la fréquence vaut $1{,}88$ pour $2025$ quand le nombre vaut
$2{,}79$, la différence venant de l'\emph{exposition}, c'est-à-dire du volume d'assurés couverts,
qui a crû de $1{,}49$. L'égalité charge $=$ nombre $\times$ coût moyen n'est donc exacte qu'avec
le nombre. La distinction sépare ce qui relève d'une croissance du portefeuille de ce qui relève
d'une dégradation du risque à portefeuille constant.

\begin{figure}[htbp]
\centering
\figover{figures_externes/lucy_taille_des_sinistres.png}
\caption[Montant indemnisé par classe de taille de sinistre]{Montant indemnisé par classe de
taille de sinistre sur le marché français, $2019$ à $2025$, en millions d'euros. Reproduit du
rapport LUCY 2026 de Nexialog Consulting \citep{RapiorKaddouri2026}. Les seuils de classe sont
ceux de la source.}
\label{fig:lucy-taille}
\end{figure}

% =====================================================================
\section{Bâle, l'ORSA et DORA : le détail du cadre}
\label{sec:ann-reglementaire}

% HARNAIS-HORS-SECTION: cadre reglementaire : numeros d'articles, delais et seuils cites des textes, donnees de droit et non resultats.

Complément du \S\ref{chap:reglementaire}.

\paragraph{Bâle III/IV : du modèle interne AMA à la standardisation SMA.}
Sous le régime de Bâle~II, les établissements bancaires d'envergure internationale
étaient encouragés à développer des modèles internes complexes, regroupés sous
l'appellation d'approche de mesure avancée (\emph{Advanced Measurement Approaches},
AMA). Ces modèles reposaient principalement sur l'approche par distribution des pertes
(\emph{Loss Distribution Approach}, LDA) ou sur des approches par scénarios. Leur flexibilité
théorique a engendré une variabilité excessive et difficilement justifiable des actifs pondérés
par le risque d'un établissement à l'autre, érodant la confiance des superviseurs dans les
ratios de solvabilité déclarés. Le présent mémoire reprend les deux marges de la LDA, fréquence
et sévérité, mais en remplace l'agrégation par une cascade dirigée
(chapitre~\ref{chap:cascade}).

La finalisation des réformes de Bâle~III, communément désignée sous le nom de Bâle~IV,
a acté le retrait des modèles internes AMA pour le calcul du capital de Pilier~1, à
compter du 1\textsuperscript{er}~janvier~2022, au profit d'une méthode unique et
comparable : l'approche d'évaluation standardisée (\emph{Standardised Measurement
Approach}, SMA). L'exigence de capital pour risque opérationnel ($ORC$) s'exprime alors
comme le produit d'une composante de volume d'activité, le \emph{Business Indicator
Component} ($BIC$), et d'un facteur d'échelle fondé sur l'historique des pertes internes,
l'\emph{Internal Loss Multiplier} ($ILM$) :
\begin{equation}
ORC = BIC \times ILM,
\qquad
RWA_{Op} = 12{,}5 \times ORC.
\end{equation}

Le \emph{Business Indicator} synthétise la taille de l'établissement à partir de trois
composantes comptables (intérêts et dividendes, services, activité financière), auxquelles
sont appliqués des coefficients marginaux croissants par tranches. Le profil de risque
propre est réintroduit par l'$ILM$, qui compare le $BIC$ à la composante de pertes
historiques $LC$, égale à quinze fois la moyenne des pertes opérationnelles annuelles
significatives des dix dernières années. Lorsque la sinistralité historique est faible
relativement à l'activité ($LC < BIC$), l'$ILM$ est inférieur à l'unité et réduit
l'exigence ; lorsqu'elle est élevée, il agit comme une pénalité en fonds propres.

\paragraph{Ce que le retrait de l'AMA ne retire pas.} Le retrait de l'AMA n'invalide pas
l'approche LDA comme outil de modélisation : il en restreint l'usage \emph{pour le calcul
réglementaire de Pilier~1 dans le secteur bancaire}. La logique de la composante $LC$, qui fait
dépendre le capital de l'historique réel des pertes, reste valide ; c'est elle, transposée à
l'assurance et enrichie de la théorie des valeurs extrêmes, que le mémoire met en oeuvre.


\paragraph{Solvabilité~II et DORA, dimension par dimension.}

\begin{center}
\renewcommand{\arraystretch}{1.3}
\begin{tabular}{>{\raggedright}p{3cm} >{\raggedright}p{5cm} >{\raggedright\arraybackslash}p{5.5cm}}
\toprule
\rowcolor{lightblue}
\textbf{Dimension} & \textbf{Solvabilité II (ORSA)} & \textbf{DORA} \\
\midrule
Taxonomie du risque & Traitement agrégé dans le module $SCR_{Op}$ &
Taxonomie granulaire des actifs TIC et indicateurs de risque distincts \\
Reporting des incidents & Notification interne, dialogue prudentiel non standardisé &
Notification obligatoire à l'ACPR : alerte sous 4\,h après classification comme majeur
(et au plus tard 24\,h après détection), rapport intermédiaire 72\,h,
rapport final 1~mois \\
Tests de résilience & Stress tests orientés adéquation du capital &
Tests de continuité et TLPT (pénétration guidée par la menace) au moins tous les 3~ans
pour les entités identifiées comme significatives \\
Cartographie des dépendances & Analyse générale des accords d'externalisation &
Registre d'Informations harmonisé, chaînes de sous-traitance jusqu'au rang~3 \\
Encadrement des tiers & Lignes directrices EIOPA non contraignantes &
Clauses contractuelles impératives (Art.~30), surveillance directe des prestataires
critiques \\
\bottomrule
\end{tabular}
\end{center}


\paragraph{Les limites structurelles de l'ORSA traditionnel.} L'ORSA
traditionnel s'appuyait principalement sur des cartographies de risques
fondées sur l'avis d'experts. Trois limites en découlaient.

\textbf{Un cloisonnement méthodologique.} Les risques TIC étaient traités séparément des
risques financiers et assurantiels, sans intégration des corrélations réelles.
L'interaction entre une défaillance de l'outil de tarification (risque opérationnel) et
la sous-évaluation systématique des primes qui en résulte (risque de souscription) était
rarement modélisée de manière dynamique.

\textbf{L'absence de données de sinistralité qualifiées.} Les assureurs ne disposaient
pas de base standardisée permettant d'alimenter les modèles de sévérité pour les
événements cyber de faible fréquence et de forte sévérité. L'étalonnage des distributions
à queue lourde restait hautement incertain et sujet à une forte volatilité de modèle,
difficulté que le mémoire traite en mobilisant plusieurs sources indépendantes
(PRC, SAS~OpRisk~Global, Hackmageddon) et en déclarant l'incertitude résiduelle.

\textbf{La non-prise en compte des interdépendances de la chaîne de valeur.} Les
politiques d'externalisation sous Solvabilité~II se limitaient à des audits documentaires
ponctuels. L'effet de cascade d'une panne systémique chez un fournisseur d'infrastructure
cloud partagé par plusieurs acteurs du marché n'était pas intégré aux tests de résistance.

\paragraph{Zone de recouvrement.} ORSA et DORA partagent un même objet, le risque opérationnel d'origine technologique,
et une même exigence d'évaluation prospective. Trois domaines se recouvrent
explicitement : le risque opérationnel informatique, traité de façon agrégée par l'ORSA
et granulaire par DORA ; les scénarios de stress cyber, que l'ORSA peut désormais alimenter
des données et tests imposés par DORA ; et la continuité d'activité, exigence commune mais
de portée différente. C'est dans ce recouvrement que se loge la contribution du présent
mémoire : utiliser les artefacts produits par DORA (taxonomie, registre, tests) pour
\emph{quantifier} dans l'ORSA ce que celui-ci ne savait jusqu'alors que décrire.

\paragraph{Zones de divergence fondamentale.}
\textbf{La finalité de l'évaluation.} Solvabilité~II s'inscrit dans une logique de
continuité \emph{financière} : prouver que l'assureur détient assez de fonds propres pour
absorber les pertes d'une défaillance majeure. DORA s'inscrit dans une logique de
continuité \emph{opérationnelle} : garantir que, même sans capital excédentaire,
l'infrastructure maintient les services critiques et se restaure rapidement. Les deux
cadres sont complémentaires, non substituables.

\textbf{Le principe de proportionnalité.} Solvabilité~II applique une proportionnalité
fondée sur la taille globale (primes, provisions). DORA décline la sienne selon la
complexité des systèmes d'information, la dépendance au cloud tiers, et la criticité dans
l'écosystème de distribution. Une captive de réassurance ou une petite mutuelle peut ainsi
bénéficier d'un régime DORA simplifié sans pour autant échapper aux exigences de capital.

\textbf{La temporalité de la supervision.} Solvabilité~II fonctionne sur des cycles
trimestriels et annuels (QRT, SFCR, RSR). DORA impose une réactivité quasi temps réel,
notamment par la notification des incidents majeurs sous quatre heures, contraignant
l'assureur à des processus permanents de veille connectés aux canaux du régulateur.

% =====================================================================
\section{Les biais de taille et de troncature : les mesures}
\label{sec:ann-biais-chiffres}

% SOURCES-SCRIPTS: 57 67

Complément du \S\ref{sec:biais-chiffres}.

\paragraph{Le biais de taille est réel, mais bien plus faible qu'annoncé.} OpRisk
sur-représente les grandes institutions, et l'on en déduit d'ordinaire que sa sévérité
sur-estime celle d'une entité moyenne. La base portant les actifs, le chiffre d'affaires et
l'effectif de la firme sinistrée, l'affirmation se teste. Une régression log-log de la perte
sur la taille donne une élasticité quasi nulle : $0{,}042$ sur les actifs
($R^2=0{,}005$), $0{,}020$ sur le chiffre d'affaires et $0{,}010$ sur l'effectif, ces deux
dernières non significatives. Prise telle quelle, cette mesure \emph{réfuterait} l'explication
usuelle.

Elle est cependant elle-même biaisée, et par le mécanisme de la section précédente. Une
petite firme n'entre dans la base que si elle subit une perte anormalement grande, quand une
grande firme y entre avec des pertes de toutes tailles : la sélection relève artificiellement
la sévérité observée des petites firmes et aplatit la pente. On corrige par un
maximum de vraisemblance sur une lognormale tronquée à gauche au seuil de collecte
($0{,}10$~M\$, $46\,\%$ des pertes de la base étant sous $1$~M\$). L'élasticité corrigée
double, $b=0{,}087$ (IC~$95\,\%$ $[0{,}026\,;0{,}148]$), et devient significative. Le seuil de
collecte était donc bien le confondant, mais l'effet reste modeste : transposée de la
firme médiane de la base ($112$~Md\$ d'actifs) à l'entité notionnelle, la sévérité est
multipliée par $0{,}85$, soit $-14{,}5\,\%$. Le biais de taille n'explique pas un facteur dix ; il
explique un septième. Ce n'est pas la sévérité qui rend le niveau non transposable,
c'est la fréquence, ce que la table~\ref{tab:echelle} établit par ailleurs. Le multiplicateur
et sa variation sont recalculés à chaque exécution, de sorte qu'ils ne peuvent plus se
désynchroniser de l'élasticité dont ils découlent.
% SOURCES-SCRIPTS: 67

\paragraph{La troncature à droite se corrige, au lieu de coûter trois années.} Les comptes
s'effondrent après 2022 ($34$ en 2023, $39$ en 2024 contre $70$ à $100$ auparavant), ce qui
conduisait à écarter ces années. La base portant la \emph{date de saisie} en plus de l'année
de survenance, on construit le triangle survenance $\times$ délai de remontée et l'on applique
un chain-ladder aux comptes. Le délai médian de remontée est de trois ans
(neuvième décile : huit ans), et les facteurs de développement successifs valent $2{,}13$,
$1{,}49$, $1{,}30$, $1{,}22$. À l'ultime, les $97$ événements observés sur 2023--2025
deviennent $255$, soit $+163\,\%$. La moyenne annuelle ultime de ces trois années
($85{,}0$) dépasse alors celle de la période pleine 2015--2022 ($74{,}1$) d'un facteur
$1{,}15$. Le risque TIC n'a pas baissé après 2022 : il n'est pas encore
remonté. Les années récentes cessent d'être une perte sèche.

\begin{figure}[htbp]\centering
\figover{S20_biais_taille_troncature.png}
\caption{Les deux biais, enfin chiffrés. (a)~la sévérité croît beaucoup moins vite
que la taille de la firme, et la pente brute est encore aplatie par le seuil de collecte.
\textbf{(b)}~les années récentes ne baissent pas, elles ne sont pas remontées : l'ultime
chain-ladder les ramène au-dessus de la moyenne historique. (c)~les deux corrections
jouent en sens opposés.}
% SOURCES-SCRIPTS: 57
\label{fig:biais}
\end{figure}

% =====================================================================
\section{Le dispositif de sources : règles du moteur et confrontation à ENISA}
\label{sec:ann-sources}

% SOURCES-SCRIPTS: 05 62 63

Complément du \S\ref{sec:portee-donnees}.

\paragraph{Deux règles du moteur de simulation.} D'une part, avec $\hat{\xi}_{\text{PRC}} =
1{,}03 \geq 1$, l'espérance de la sévérité PRC est infinie et la VaR simulée
diverge numériquement ; un plafond de sévérité de 40~M\euro, ancré sur la
capacité de réassurance cyber observable sur le marché, est appliqué à la
\emph{seule source PRC}. La calibration OpRisk, dont l'indice de queue
$\hat{\xi}_{\text{OpRisk}} = 0{,}595 < 1$ garantit une espérance finie, n'est
en revanche pas plafonnée : ses montants étant observés de bout en bout, y
appliquer un plafond tronquerait des pertes réelles et contredirait sa propre
calibration. D'autre part, la sévérité de queue n'est tirée que pour la
fraction $p_u$ des événements qui dépassent le seuil ; les autres reçoivent une
sévérité nulle, ce qui fait du besoin publié un capital de queue.

La conséquence directe de ce cadre épistémique est méthodologique autant que
présentationnelle : le résultat final ne peut être livré comme un montant
unique, mais doit être présenté comme une distribution large, conditionnelle
à des hypothèses documentées et à des sources dont les écarts sont eux-mêmes
informatifs, et ce principe structure tout le chapitre~\ref{chap:resultats}.

%--------------------------------------------------------------

%--------------------------------------------------------------

Le dispositif retenu (OpRisk Global pour la sévérité observée et la fréquence
d'entrée, PRC pour la sévérité alternative, Hackmageddon pour la structure du
risque) appelle une
question préalable à toute discussion de robustesse : ce triptyque est-il
suffisant pour l'exercice mené dans ce mémoire ? La réponse tient à la nature
même de l'exercice. L'objectif n'est pas de produire une formule
réglementaire de Pilier~1, qui exigerait une base exhaustive et représentative
du marché au sens de l'Article~19, base qui, comme établi plus haut, n'existe aujourd'hui
pour aucun acteur, public ou privé, sur le
risque cyber. L'objectif est de construire, dans l'esprit du Pilier~2 et de
l'ORSA, une mesure prospective dont la valeur ajoutée actuarielle réside dans
la transparence de ses limites autant que dans la précision de son résultat.
Sous cet angle, les trois sources mobilisées couvrent exactement les trois
rôles qu'exige une architecture fréquence-sévérité-structure, et leurs biais
respectifs (sélection, déclaration, survie, visibilité, instabilité
taxonomique) sont documentés plutôt que dissimulés.


La structure du risque tirée de Hackmageddon (\S\ref{sec:hackmageddon}), bien que
documentée quant à ses limites de visibilité et de taxonomie, gagne à être
confrontée à une source indépendante, de nature différente et couvrant un
périmètre réglementaire proche de celui du présent mémoire. Le rapport
\emph{ENISA Threat Landscape 2025} \citep{enisa2025}, qui ancre au \S\ref{sec:quatre-canaux} les
bornes du multiplicateur associé à l'exploitation de vulnérabilités, remplit
ce rôle : constitué par l'agence européenne de cybersécurité elle-même à
partir de 4\,875 incidents de périmètre européen, il repose sur une
méthodologie de collecte distincte de celle, essentiellement médiatique, de
Hackmageddon.

La confrontation porte sur un point précis, celui qui ancre cette fourchette : ENISA
mesure un taux de conversion tentative-vers-intrusion de 70~\% pour
l'exploitation de vulnérabilités, contre seulement 27~\% pour le phishing.
Cette asymétrie, mesurée par un organisme réglementaire sur un périmètre
indépendant, confirme un résultat qu'Hackmageddon, par construction, ne peut
révéler à partir du seul volume d'incidents observés : le phishing domine la
structure en volume (38{,}8~\% contre 33{,}8~\% pour l'exploitation de
vulnérabilités), mais convertit nettement moins efficacement en
intrusion réussie. Les deux sources s'accordent ainsi, chacune sur son propre
plan (structure des incidents rapportés pour Hackmageddon, taux de
conversion pour ENISA) sur une même hiérarchie qualitative des vecteurs
d'attaque, ce qui renforce la crédibilité du mapping vecteurs~$\rightarrow$~exigences
DORA retenu au chapitre~\ref{chap:multietats} sans nécessiter de recalibration.

\paragraph{Piste non mobilisée : Verizon DBIR.} Le \emph{Data Breach
Investigations Report} de Verizon, fondé sur des méthodologies de collecte
distinctes à la fois de Hackmageddon et d'ENISA, constituerait une troisième
confrontation naturelle. Il n'est pas mobilisé dans le présent mémoire, faute
de statistiques suffisamment vérifiées dans le périmètre de ce travail pour
en garantir l'exactitude des chiffres cités. Cette confrontation constitue une
perspective directe de prolongement, dans le même esprit que l'automatisation
du reclassement par NLP évoquée plus haut.

% =====================================================================
\section{Les théorèmes des valeurs extrêmes, énoncés développés}
\label{sec:ann-evt}

% HARNAIS-HORS-SECTION: enonces de theoremes : constantes d'enonce, aucune valeur calculee par le pipeline.

Là où le Théorème Central Limite s'intéresse au comportement asymptotique des sommes de variables aléatoires, la Théorie des Valeurs Extrêmes (EVT) étudie le comportement asymptotique de leurs maxima.

Soit une séquence $(X_1, X_2, \dots, X_n)$ de variables aléatoires indépendantes et identiquement distribuées (i.i.d.) de fonction de répartition $F$, représentant la sévérité des cyberattaques. L'approche initiale de l'EVT, dite des \textit{Block Maxima}, s'intéresse au maximum sur des blocs de taille $n$ (par exemple, le sinistre cyber maximal annuel) :
\begin{equation}
M_n = \max(X_1, X_2, \dots, X_n)
\end{equation}



\paragraph{Théorème de Fisher-Tippett-Gnedenko (1928, 1943).} S'il existe des suites de constantes de normalisation $a_n > 0$ et $b_n \in \mathbb{R}$, et une distribution non dégénérée $H$ telles que :
\begin{equation}
\lim_{n \to \infty} \mathbb{P}\left( \frac{M_n - b_n}{a_n} \le x \right) = \lim_{n \to \infty} F^n(a_n x + b_n) = H(x)
\end{equation}
Alors, $H$ appartient nécessairement à la famille des distributions des extrêmes généralisées (GEV - \textit{Generalized Extreme Value}), dont la fonction de répartition s'écrit :
\begin{equation}
H_{\xi}(x) =
\begin{cases}
\exp\left(-\left(1 + \xi x \right)^{-1/\xi}\right) & \text{si } \xi \neq 0 \\
\exp\left(-\exp(-x)\right) & \text{si } \xi = 0
\end{cases}
\end{equation}
définie sur l'ensemble $\{x : 1 + \xi x > 0\}$. Le paramètre $\xi$ (indice de queue) définit trois grands domaines d'attraction :
\begin{itemize}
    \item \textbf{$\xi > 0$ (Loi de Fréchet) :} La distribution présente une queue lourde à décroissance polynomiale. C'est le régime qui décrit les sinistres cyber les plus lourds.
    \item \textbf{$\xi = 0$ (Loi de Gumbel) :} La distribution présente une queue légère à décroissance exponentielle (domaine d'attraction des lois Normale, Lognormale, Gamma).
    \item \textbf{$\xi < 0$ (Loi de Weibull) :} La distribution a un support fini supérieur (queue courte).
\end{itemize}






%--------------------------------------------------------------

%--------------------------------------------------------------

L'approche \textit{Block Maxima} (GEV) est rigoureuse mais coûteuse en information. En ne
retenant que le maximum de chaque bloc, elle écarte le deuxième plus grand sinistre d'une année
extrême, lequel peut dépasser le maximum d'une année calme. Sur une donnée cyber aussi peu
fournie, ce rejet d'observations fragiliserait l'ajustement.

Pour contourner ce problème, ce mémoire privilégie l'approche \textit{Peaks-Over-Threshold} (POT). Au lieu de ne garder que les maxima, cette méthode exploite toutes les observations excédant un seuil élevé $u$. On s'intéresse alors à la distribution des excès $Y = X - u$.

La fonction de répartition conditionnelle des excès est donnée par :
\begin{equation}
F_u(y) = \mathbb{P}(X - u \le y \mid X > u) = \frac{F(u+y) - F(u)}{1 - F(u)}
\end{equation}

\paragraph{Théorème de Balkema-de Haan-Pickands (1974).} Ce théorème établit le pont mathématique entre la GEV et l'approche POT. Pour une large classe de distributions sous-jacentes $F$ (celles appartenant au domaine d'attraction d'une loi GEV $H_\xi$), il existe une fonction mesurable positive $\sigma(u)$ telle que, lorsque le seuil $u$ tend vers la limite supérieure du support de $F$ ($x_F$), la distribution des excès converge asymptotiquement vers une distribution de Pareto Généralisée (GPD) :
\begin{equation}
\lim_{u \to x_F} \sup_{0 \le y < x_F - u} \left| F_u(y) - G_{\xi, \sigma(u)}(y) \right| = 0
\end{equation}

Où la fonction de répartition de la GPD s'écrit :
\begin{equation}
G_{\xi, \sigma}(y) =
\begin{cases}
1 - \left(1 + \frac{\xi y}{\sigma}\right)^{-1/\xi} & \text{si } \xi \neq 0 \\
1 - \exp\left(-\frac{y}{\sigma}\right) & \text{si } \xi = 0
\end{cases}
\end{equation}
avec $y \ge 0$ si $\xi \ge 0$, et $0 \le y \le -\sigma/\xi$ si $\xi < 0$. Le paramètre de forme $\xi$ de la GPD est strictement le même que celui de la loi GEV correspondante. Dans le contexte du risque cyber, un indice de queue $\xi > 0$ traduit une distribution à queue lourde de type Fréchet, impliquant des moments potentiellement infinis (l'espérance est infinie si $\xi \ge 1$, la variance si $\xi \ge 0.5$).





%--------------------------------------------------------------

%--------------------------------------------------------------



\begin{proof}[Éléments pour le théorème~\ref{soc:th:ftg}]
La loi limite $H$ est nécessairement \emph{max-stable} : pour tout $n$, $H^n(a_nx+b_n)=H(x)$, car $M_{nk}$ est le maximum de $k$ blocs de taille $n$. L'équation fonctionnelle de max-stabilité, résolue via la théorie des fonctions à variation régulière, n'admet (à normalisation près) que la famille à un paramètre $H_\xi$. La correspondance queue de $F$ $\leftrightarrow$ signe de $\xi$ s'obtient en caractérisant les $F$ dont la fonction de survie est à variation régulière d'indice $-1/\xi$ (cas Fréchet). Le détail figure en Annexe~\ref{soc:ann:ftg}.
\end{proof}

\begin{proof}[Esquisse pour le théorème~\ref{soc:th:balkema}]
L'appartenance de $F$ au domaine d'attraction de $H_\xi$ équivaut à une condition de variation régulière sur la fonction de survie $\bar F$. On montre que $\bar F(u+y\,a(u))/\bar F(u)\to(1+\xi y)^{-1/\xi}$ quand $u\to x_F$, pour une fonction d'échelle $a(u)$ adéquate ; ce membre de droite est exactement la fonction de survie GPD. La convergence uniforme découle de la monotonie des fonctions de répartition (théorème de Pólya). Détail et lien avec le Théorème~\ref{soc:th:ftg} en Annexe~\ref{soc:ann:balkema}.
\end{proof}

% =====================================================================
\section{L'estimation de l'indice de queue : lois asymptotiques et démonstration}
\label{sec:ann-estimation-queue}

% SOURCES-SCRIPTS: 47 63 67

\paragraph{L'estimateur de Hill.}
Pour $\xi > 0$, l'estimateur non paramétrique de Hill (1975) estime l'indice de queue à partir des $k$ plus grandes statistiques d'ordre, sans passer par l'ajustement d'une loi paramétrique. Asymptotiquement convergent et sans biais lorsque la loi appartient au domaine d'attraction de Fréchet, il présente en échantillon fini un biais majeur : si la composante à variation lente de la queue ne converge pas assez vite, $\hat{\xi}_{\text{Hill}}$ devient instable.

\paragraph{Loi asymptotique et information de Fisher.} La qualité d'un intervalle de confiance sur $\xi$ conditionne toute la lecture du capital. Le fondement théorique en est donné ici ; il sert de contrôle indépendant du bootstrap employé au chapitre~\ref{chap:resultats}.

\begin{proposition}[Normalité asymptotique de l'EMV de la GPD]\label{soc:prop:fisher}
Pour $\xi > -\tfrac{1}{2}$, l'estimateur du maximum de vraisemblance $(\hat\xi,\hat\sigma)$ de la GPD est consistant et asymptotiquement normal :
\begin{equation}
\sqrt{N_u}\,\begin{pmatrix}\hat\xi-\xi\\[2pt]\hat\sigma-\sigma\end{pmatrix}
\;\xrightarrow{\;d\;}\;
\mathcal{N}\!\left(\mathbf{0},\; (1+\xi)\begin{pmatrix}1+\xi & -\sigma\\[2pt] -\sigma & 2\sigma^2\end{pmatrix}\right).
\end{equation}
En particulier, la variance asymptotique de l'indice de queue est
$\Var(\hat\xi)\simeq (1+\xi)^2/N_u$.
\end{proposition}

\begin{proof}[Esquisse]
La matrice de covariance limite est l'inverse de l'information de Fisher $\mathcal{I}(\xi,\sigma)$, obtenue en prenant l'espérance de l'opposé de la hessienne de $\ell$. Le détail du calcul de $\mathcal{I}$ (dérivées secondes de \eqref{soc:eq:loglik-gpd} et intégration contre la densité GPD, valide pour $\xi>-\tfrac12$ où l'information est finie) est reporté en Annexe~\ref{soc:ann:fisher}. Le résultat classique remonte à \citet{smith1987}.
\end{proof}

\noindent Appliquée à la calibration OpRisk ($\hat\xi=0{,}595$, $N_u=91$ excès), la Proposition~\ref{soc:prop:fisher} donne un écart-type asymptotique $\widehat{\mathrm{sd}}(\hat\xi)=(1+\hat\xi)/\sqrt{N_u}=0{,}167$, soit un intervalle de confiance à 90\,\% par la delta-méthode de $[0{,}320\,;\,0{,}870]$. Cet intervalle recoupe presque exactement l'intervalle bootstrap $[0{,}304\,;\,0{,}831]$ obtenu indépendamment au \S\ref{soc:sec:severite} : la concordance des deux approches, l'une paramétrique et asymptotique, l'autre par rééchantillonnage, renforce la fiabilité de l'incertitude reportée sur~$\xi$.

\begin{proposition}[Loi asymptotique de l'estimateur de Hill]\label{soc:prop:hill}
Si $F$ est à variation régulière d'indice $-1/\xi$ (domaine de Fréchet), alors pour une suite $k=k(n)\to\infty$ avec $k/n\to0$,
\begin{equation}
\sqrt{k}\,\bigl(\hat\xi_{Hill}(k)-\xi\bigr)\;\xrightarrow{\;d\;}\;\mathcal{N}(0,\;\xi^2).
\end{equation}
\end{proposition}

\noindent L'écart-type $\xi/\sqrt{k}$ décroît en $\sqrt{k}$ : d'où l'arbitrage biais--variance qui gouverne le choix de $k$ dans le \emph{Hill plot}, et qui explique l'instabilité signalée plus haut lorsque $k$ est pris trop grand (contamination par le corps de la loi) ou trop petit (variance explosive). La preuve, fondée sur la représentation de Rényi des espacements exponentiels, est esquissée en Annexe~\ref{soc:ann:hill}.

\noindent Démonstration de la proposition~\ref{soc:prop:varestvar}.

\begin{proof}
Pour $x>u$, la propriété POT donne $\Prob(X>x)=\zeta_u\bigl(1+\xi(x-u)/\sigma\bigr)^{-1/\xi}$. En posant ce membre égal à $1-\alpha$ et en inversant, on obtient \eqref{soc:eq:var-gpd}. Pour \eqref{soc:eq:tvar-gpd}, la stabilité de la GPD par franchissement de seuil implique que l'excès au-delà de $\qsev_\alpha(X)$ suit encore une GPD de paramètre $\xi$ et d'échelle $\sigma+\xi(\qsev_\alpha(X)-u)$, d'espérance $\bigl(\sigma+\xi(\qsev_\alpha(X)-u)\bigr)/(1-\xi)$ pour $\xi<1$ ; en ajoutant $\qsev_\alpha(X)$ et en réarrangeant, on retrouve \eqref{soc:eq:tvar-gpd}. La divergence pour $\xi\ge1$ découle de $\E[X]=\int_u^\infty \Prob(X>x)\,\dd x=+\infty$ dès que $1/\xi\le1$. Détail en Annexe~\ref{soc:ann:varestvar}.
\end{proof}

\noindent La Proposition~\ref{soc:prop:varestvar} éclaire par ailleurs le contraste entre sources : la calibration PRC affiche $\hat\xi=1{,}033>1$, ce qui place la sévérité en régime d'espérance infinie : la moyenne de queue, et donc tout Expected Shortfall bâti sur elle, y est mathématiquement indéfinie. C'est ce qui justifie \emph{a posteriori} le recours à un plafond de sévérité (capacité de réassurance) pour rendre le capital PRC calculable, et c'est aussi le premier argument contre l'usage d'une mesure de queue comme \emph{niveau de couverture} : sur l'une des deux sources, cette mesure n'existe pas.

% =====================================================================
\section{Hill contre maximum de vraisemblance : le détail}
\label{sec:ann-hill}

% SOURCES-SCRIPTS: 47 63 67

Complément du \S\ref{soc:sec:severite}.

\paragraph{L'estimateur de Hill : définition et rôle diagnostique.}
Introduit par \citet{Hill1975}, l'estimateur de Hill fournit une estimation
de l'indice de queue $\xi$ à partir des $k$ statistiques d'ordre les plus
élevées d'un échantillon $X_{(1)} \leq \dots \leq X_{(n)}$, sans passer par
l'ajustement paramétrique d'une loi de Pareto généralisée :
\begin{equation}
\hat{\xi}_{\text{Hill}}(k) = \frac{1}{k}\sum_{i=1}^{k}
\ln X_{(n-i+1)} - \ln X_{(n-k)}.
\label{soc:eq:hill}
\end{equation}
Contrairement au maximum de vraisemblance sur la GPD, qui exploite le cadre
théorique de Pickands--Balkema--de~Haan et estime conjointement $\xi$ et
$\sigma$, l'estimateur de Hill repose sur une hypothèse de Pareto \emph{pure}
dans la queue et ne fournit qu'un indice de forme, sans paramètre d'échelle
explicite \citep{Beirlant2004}. Sa propriété la plus documentée dans la
littérature EVT est un compromis biais--variance piloté par le choix de
$k$~: pour $k$ trop petit, l'estimateur est instable (variance élevée,
intervalle de confiance large) ; pour $k$ trop grand, il intègre des
observations situées hors du régime purement asymptotique de la queue, ce
qui induit un biais positif croissant avec $k$ \citep{Embrechts1997,
Beirlant2004}.

\paragraph{Application et confrontation au MLE-GPD.} Le Hill plot obtenu sur
le périmètre cyber$\times$finance d'OpRisk (figure~\ref{soc:fig:hill-ic})
illustre ce compromis de façon nette~: l'estimateur croît de façon quasi
monotone avec $k$, passant de $\hat{\xi} \approx 0{,}5$ pour $k$ faible à
$\hat{\xi} \approx 1{,}8$ pour $k = 200$, tandis que la bande de confiance
asymptotique à 90~\% s'élargit continûment. La référence par maximum de
vraisemblance ($\hat{\xi}_{\text{MLE}} = 0{,}595$, pointillés rouges) n'est
statistiquement compatible avec l'estimateur de Hill que sur une plage
restreinte, $k \in [3, 25]$ environ, où la bande de confiance recouvre
encore la valeur MLE. Au-delà de $k \approx 30$, l'estimateur de Hill
diverge continûment de cette référence, ce qui traduit le biais positif
attendu lorsque des observations hors du régime purement asymptotique de la
queue sont intégrées au calcul. Sur le seuil effectivement retenu pour la
calibration ($k = 86$), cet écart atteint 120{,}8~\%
($\hat{\xi}_{\text{Hill}} = 1{,}32$ contre $\hat{\xi}_{\text{MLE}} = 0{,}595$),
ce qui confirme que l'estimateur de Hill n'est pas adapté à ce niveau de
$k$ et que le MLE-GPD, qui exploite un cadre théorique plus général
(loi GPD complète plutôt que Pareto pure), demeure l'estimateur de
référence pour la calibration finale.

\paragraph{Cet écart est démontré, pas invoqué, et l'enjeu justifie qu'on s'y arrête.}
Écrire que l'estimateur de Hill est biaisé en échantillon fini est vrai en général, mais c'est
un argument d'autorité : cela ne dit pas si \emph{cet} écart-là, sur \emph{ces} données, relève
du biais ou d'un désaccord réel. Or l'enjeu n'est pas mince. À $\hat\xi = 0{,}595$, la variance
de la sévérité est infinie mais l'espérance existe ; à $\hat\xi = 1{,}32$, l'espérance elle-même
cesse d'exister, et le modèle de perte agrégée qui suit n'aurait plus d'objet. C'est le point le
plus attaquable du socle, et il se tranche par le calcul plutôt que par la citation.

On simule donc $2\,000$ échantillons sous le modèle \emph{publié}, c'est-à-dire le corps observé
sous le seuil et une queue GPD de paramètres $(\hat\xi,\hat\sigma) = (0{,}595\,;57{,}97)$
au-dessus, puis on applique à chacun exactement le même estimateur de Hill, aux mêmes $k$. Si un
monde où $\xi$ vaut $0{,}595$ reproduit les valeurs de Hill observées, alors ces valeurs ne
témoignent de rien : elles sont précisément ce qu'un tel $\xi$ produit.

\begin{center}\footnotesize
\renewcommand{\arraystretch}{1.25}
\begin{tabular}{@{}r r r c@{}}
\toprule
$k$ & $\hat\xi_{\text{Hill}}$ observé & simulé (moyenne) & intervalle simulé à $90\,\%$ \\
\midrule
$30$  & $0{,}932$ & $0{,}825$ & $[0{,}604\,;1{,}076]$ \\
$50$  & $0{,}956$ & $0{,}961$ & $[0{,}768\,;1{,}165]$ \\
$86$  & $1{,}315$ & $1{,}327$ & $[1{,}155\,;1{,}506]$ \\
$91$  & $1{,}396$ & $1{,}422$ & $[1{,}248\,;1{,}596]$ \\
$120$ & $1{,}457$ & $1{,}477$ & $[1{,}344\,;1{,}608]$ \\
$200$ & $1{,}717$ & $1{,}728$ & $[1{,}649\,;1{,}807]$ \\
\bottomrule
\end{tabular}
\end{center}

\vspace{4pt}
\noindent Les six valeurs observées tombent dans l'intervalle simulé, et les moyennes coïncident
à quelques millièmes près. L'écart de $120{,}8\,\%$ n'est donc pas un désaccord entre deux
estimateurs : c'est le comportement attendu de Hill appliqué à des pertes brutes, décalées du
seuil, sur un échantillon de cette taille. Les données ne contredisent pas la calibration
publiée, elles la corroborent par un troisième chemin, après Anderson-Darling et
Kolmogorov-Smirnov.

Deux dérives de sens contraire achèvent de le montrer. Un indice de queue authentique produit un
\emph{plateau} sur le tracé de Hill ; ici l'estimateur croît sans s'arrêter avec $k$, de
$0{,}93$ à $k=30$ jusqu'à $1{,}72$ à $k=200$, c'est-à-dire qu'il monte à mesure qu'on descend
dans le corps de la loi. Le balayage de seuil va exactement dans l'autre sens, puisque $\hat\xi$
\emph{décroît} de $0{,}60$ à $0{,}38$ quand on remonte dans la queue. Si la queue valait
réellement $1{,}32$, les deux dérives seraient inversées.
% SOURCES-SCRIPTS: 47

\begin{figure}[htbp]
\centering
\includegraphics[width=0.85\textwidth]{figures/hill_plot_ic90_oprisk.png}
\caption{Hill plot avec intervalle de confiance asymptotique à 90~\% sur
OpRisk Global. La référence MLE ($\hat{\xi} \approx 0{,}595$, pointillés
rouges) n'est statistiquement compatible avec l'estimateur de Hill que pour
$k \lesssim 25$~; au-delà, la dérive de $\hat{\xi}_{\text{Hill}}$ traduit le
biais croissant lié à l'intégration d'observations hors du régime
asymptotique pur de la queue.}
\label{soc:fig:hill-ic}
\end{figure}



\paragraph{Validation croisée PRC.} La calibration sur la PRC, via la
conversion Jacobs, donne $\hat{\xi}_{\text{PRC}} = 1{,}03$ (seuil
$u = 4{,}18$~M\euro, 2\,258 excès), supérieur à
$\hat{\xi}_{\text{OpRisk}} = 0{,}595$, les IC90\% des deux estimateurs ne se
chevauchant pas. Ce paramètre correspond formellement à une espérance de
sévérité infinie ($\xi \geq 1$), situation que la littérature du risque opérationnel
traite explicitement \citep{neslehova2006}, ce qui impose un plafond de sévérité
opérationnel, fixé à 40~M\euro{} et ancré sur la capacité de réassurance cyber
observable sur le marché, pour que la simulation reste numériquement
stable ; ce plafond ne s'applique qu'à la source PRC, la calibration OpRisk
($\xi < 1$, espérance finie) demeurant non plafonnée. L'écart entre
les deux calibrations ne mesure pas une différence de risque : la sévérité PRC
est une puissance du nombre d'enregistrements, et une puissance multiplie
l'indice de queue par son exposant, si bien que les deux indices ne portent pas
sur le même objet (\S\ref{sec:oprisk}). Les deux sources confirment
néanmoins le même régime qualitatif de queue lourde ($\xi > 0$), et aucune
hiérarchie de type \og source primaire\fg{} n'est retenue entre elles : leur
écart est traité comme une composante de l'incertitude totale du modèle,
non comme un désaccord à trancher.

% =====================================================================
\section{Le seuil de la calibration : deux règles écrites}
\label{sec:ann-seuil}

% SOURCES-SCRIPTS: 93 47

Complément du \S\ref{soc:sec:validation}.

\paragraph{Le seuil est-il un choix ou une règle ?}
Tout ce qui précède décrit le comportement de l'ajustement autour du seuil retenu. Un lecteur
demande autre chose : pourquoi \emph{ce} seuil, et qu'aurait choisi une règle écrite. Deux règles
sont donc mises en \oe uvre sur dix percentiles candidats, du $70^{\text{e}}$ au $95^{\text{e}}$.
La première retient le plus \emph{bas} seuil que le test d'adéquation ne rejette pas, puisque
descendre gagne des excès ; la seconde retient le plus bas seuil au-dessus duquel l'indice de
queue cesse de dériver, c'est-à-dire tel que tous les indices estimés plus haut tombent dans son
intervalle de confiance à $90\,\%$.

\textbf{La règle de stabilité sélectionne exactement le seuil publié}, $19{,}87$~M\euro{} contre
$20{,}03$, soit un écart de $0{,}8\,\%$ et un quantile identique au centième de pour cent. Le
seuil publié est donc précisément le plus bas au-dessus duquel $\hat\xi$ se stabilise, ce qui
n'est pas un critère fabriqué après coup : c'est le balayage décrit plus haut, employé comme
règle de sélection au lieu d'être seulement rapporté. Et c'est aussi le seuil le \emph{mieux
ajusté} du balayage, sa $p$-value d'Anderson-Darling valant $0{,}971$ quand le second meilleur
candidat vaut $0{,}954$ : deux critères indépendants désignent le même point, ce qui n'était pas
garanti.

\textbf{La règle la plus permissive descendrait plus bas et donnerait davantage de capital}, et
c'est l'argument le plus solide contre le soupçon d'un choix opportuniste. Elle retient
$12{,}88$~M\euro, soit $35{,}7\,\%$ sous le seuil publié, où le quantile de sévérité vaut
$758{,}9$ contre $654{,}5$, soit $16{,}0\,\%$ de plus. Le seuil publié n'est donc pas celui qui
maximise le chiffre du mémoire : parmi les candidats acceptables, il est du côté qui le minore.

\begin{attention}
\textbf{Et le compromis biais-variance que l'on écrit par habitude n'existe pas sur cette
grandeur.} On attend qu'un seuil plus bas gagne en précision ce qu'il perd en biais, puisqu'il
ajoute des excès. La mesure dit l'inverse : l'écart-type du quantile à $99{,}5\,\%$ passe de
$519{,}8$ à $150{,}3$~M\euro{} en remontant le seuil, soit une division par $3{,}5$, alors que le
nombre d'excès est divisé par presque six, de $174$ à $30$. Le motif est propre à la grandeur
visée : le quantile dépend \emph{exponentiellement} de l'indice de queue, et descendre le seuil
gonfle cet indice jusqu'à $1{,}058$ au percentile $70$, donc au-delà de l'unité où l'espérance de
la sévérité cesse d'exister. Les excès gagnés en descendant sont des pertes qui n'appartiennent
pas encore à la queue : elles ajoutent du biais, et ce biais se propage à l'écart-type par le
levier de l'indice. Il n'y a donc rien à arbitrer entre biais et précision, les deux allant du
même côté, et le seuil doit se choisir sur l'adéquation et la stabilité. Un dernier avertissement
de lecture : le faible écart-type du percentile $95$ ne se lit pas comme de la précision, étant
estimé sur $30$ excès seulement, donc lui-même mal connu.
\end{attention}

% =====================================================================
\section{Le backtest hors échantillon : le détail des tests}
\label{sec:ann-backtest}

% SOURCES-SCRIPTS: 89 90 91 102

Complément du \S\ref{soc:sec:backtest}.

\paragraph{La loi de fréquence est validée, et c'est le résultat le plus net.} Les comptes
annuels du périmètre ont pour moyenne $24{,}50$ et pour variance $58{,}55$, soit un indice de
dispersion de $2{,}39$, sans tendance significative ($-0{,}092$ incident par an, $p=0{,}7297$) :
une loi sans mémoire n'est donc pas mal spécifiée sur ce point. Sur les douze années notées, un
intervalle de prévision annoncé à $90\,\%$ en couvre $91{,}7\,\%$ sous la binomiale négative
contre $75{,}0\,\%$ sous le Poisson. La couverture ne suffit pourtant pas à trancher, un
intervalle inutilement large pouvant couvrir sans rien valoir. Le log-score prédictif
cumulé, qui pénalise la largeur autant que l'étroitesse, désigne le même gagnant :
$-41{,}334$ contre $-44{,}388$, soit $3{,}054$ nats sur douze années. Le paramètre
supplémentaire de la binomiale négative est donc gagné hors échantillon, et la surdispersion
cesse d'être une précaution de modélisation pour devenir une propriété mesurée deux fois, dans
l'échantillon par le rapport de vraisemblance et hors échantillon par la prévision.

\paragraph{La forme de la queue survit, et ce n'était pas acquis.} Chaque ajustement d'origine ne
dispose que de vingt à quatre-vingt-dix excès. La transformée de probabilité intégrale des excès
de test, qui doit être uniforme si la loi ajustée est la bonne, a pour moyenne $0{,}532$ et n'est
pas rejetée par le test de Kolmogorov-Smirnov ($D=0{,}100$, $p=0{,}4087$). L'écart de cette
moyenne à $0{,}500$ va dans le sens de pertes tombant plus haut dans la queue ajustée que prévu,
mais il n'est \emph{pas} significatif, et l'énoncer autrement serait surinterpréter un test qui
ne rejette pas. Ce qui est établi est plus modeste et suffit : la GPD n'est pas la mauvaise
famille.

\paragraph{Les quantiles, en revanche, sont dépassés trois fois trop souvent.} Sur les $280$
pertes notées, le quantile de sévérité que l'apprentissage annonçait à $99\,\%$ est franchi huit
fois pour $2{,}8$ attendues ($p=0{,}0078$), et celui à $99{,}5\,\%$ cinq fois pour $1{,}4$
($p=0{,}0140$). Le test de dépassement rejette donc aux deux niveaux là où le test de forme ne
rejetait pas. Les deux constats semblent s'opposer et ils ne portent pas sur le même objet : la
transformée intégrale ne regarde que les excès au-dessus du seuil, donc la queue
\emph{conditionnelle}, quand un quantile combine la forme, le seuil \emph{et} le taux de
dépassement. Si la forme tient et le quantile non, c'est le taux qui dérive, et il dérive :
$15{,}1\,\%$ en apprentissage par construction de la règle, $27{,}1\,\%$ hors échantillon, un
rapport de $1{,}80$ que le test binomial situe à $p=2{,}35\cdot 10^{-7}$. La dérive se lit
d'ailleurs directement sur le seuil, qui est un quantile de l'apprentissage et passe de
$8{,}56$ à $18{,}20$~M\euro{} en onze ans, ainsi que sur la médiane annuelle des pertes, qui
monte de $13{,}0\,\%$ par an en logarithme ($p=0{,}0003$). Le défaut est donc un défaut de
stationnarité de l'échelle, non un défaut de famille.

\paragraph{Ce que cette dérive coûte, chiffré plutôt que corrigé.} Mesurer une dérive ne dit pas
ce qu'elle déplace sur le nombre publié. Le calcul porte sur la donnée complète de la chaîne
publiée, $582$ incidents dont $91$ excès au-dessus du seuil de $20{,}03$~M\euro. Un ajustement
stationnaire y donne un quantile de sévérité à $99{,}5\,\%$ de $657{,}16$~M\euro, à comparer aux
$662{,}78$ publiés : l'écart entre ces deux valeurs est le défaut de taux de dépassement déjà
déclaré au \S\ref{sec:pu-coherence}, et les variantes qui suivent sont donc rapportées à la
première plutôt qu'à la seconde, sans quoi deux écarts se mélangeraient dans un seul nombre.
L'échelle est ensuite rendue non stationnaire, $\sigma_t = \sigma\,e^{\,b\,(t-t_{\mathrm{ref}})}$
à indice de queue constant, et estimée par maximum de vraisemblance sur les mêmes excès. La
dérive de la queue vaut alors $4{,}00\,\%$ par an, d'écart-type $1{,}84\,\%$, et le rapport de
vraisemblance la déclare significative ($4{,}28$, $p=0{,}0386$). Le quantile évalué en fin de
période passe à $819{,}40$~M\euro, soit $162$~M\euro{} et $24{,}7\,\%$ de plus. Cet écart est une
borne basse : le taux de dépassement est tenu commun aux deux branches, alors qu'une
échelle qui monte fait aussi monter le nombre d'excès au-dessus d'un seuil fixe.

\paragraph{Mais la dérive n'est pas la même selon l'endroit de la distribution, ce qui modère le
paragraphe précédent.} Le corps dérive de $13{,}0\,\%$ par an quand la queue ne dérive que de
$4{,}00\,\%$, un facteur $3{,}2$. Lire la dérive du corps comme une dérive du capital serait donc
une erreur d'un facteur trois, et c'est précisément la lecture qu'invite le paragraphe précédent
pris seul. La contre-épreuve le confirme par l'absurde : indexer toutes les pertes à la tendance
de la médiane, puis rejouer la chaîne à percentile de seuil constant, donne un quantile de
$4\,291$~M\euro, soit $6{,}5$ fois la référence, et surtout un indice de queue de $1{,}27$ à la
borne haute de l'intervalle de la tendance. Au-delà de l'unité l'espérance de la sévérité cesse
d'exister, et une hypothèse qui détruit l'objet qu'elle prétend corriger ne se retient pas.
C'est le même argument d'existence que celui qui écarte la source PRC comme support d'une mesure
de couverture, et il vaut mieux qu'un argument de degré.

\paragraph{Une part de la dérive est monétaire.} Les montants de la chaîne publiée sont
nominaux (chapitre~\ref{chap:donnees}). Rejouée sur la valeur courante de la même base, qui les
corrige d'une inflation implicite de $2{,}19\,\%$ par an, la tendance de la médiane passe de
$13{,}0$ à $10{,}6\,\%$ par an et reste significative ($p=0{,}0021$) : l'inflation générale n'en
explique que $17\,\%$. Dans la queue, au même seuil publié désormais franchi par $106$ excès, la
dérive tombe de $4{,}00$ à $2{,}00\,\%$ par an et cesse d'être significative ($p=0{,}211$). La
moitié de la dérive de queue est donc de l'inflation monétaire, et le reste ne se distingue pas
du bruit sur cet historique. Pour un capital exprimé en euros courants, le sens de l'écart
déclaré ne change pas : une perte ancienne comptée à sa valeur nominale est sous-évaluée,
qu'elle le soit par l'inflation ou par une aggravation du risque. Ce qui change est sa lecture :
la hausse de l'échelle réelle des pertes de queue n'est pas démontrée.
% SOURCES-SCRIPTS: 102

\paragraph{Une conséquence sur l'indice de queue publié, et elle va en sens contraire.}
Modéliser la dérive fait \emph{tomber} l'indice de queue, de $0{,}5979$ à $0{,}5273$. Le
mécanisme n'a rien de subtil : regrouper des excès d'années dont les échelles diffèrent fabrique
un mélange, et un mélange de lois à échelles inégales paraît plus lourd de queue qu'aucune de ses
composantes. L'ajustement stationnaire attribue donc à la \emph{forme} une part de ce qui relève
de la \emph{dérive}. Les deux écarts du modèle publié sont ainsi de sens contraires, l'indice de
queue étant légèrement prudent quand l'échelle est anti-conservatrice, et ils ne se compensent
pas puisqu'ils ne portent pas sur la même grandeur.

\paragraph{L'effet sur l'écart entre états est mesuré, et l'argument attendu ne vaut que pour
l'échelle.} Il serait naturel d'affirmer qu'une dérive commune aux deux états de conformité se
simplifie dans un rapport, donc que l'écart DORA n'est pas concerné. L'affirmation est exacte
pour une variation de la seule \emph{échelle} de sévérité, et elle est désormais mesurée plutôt
qu'avancée : multiplier l'échelle multiplie les deux états par le même facteur, $1{,}4556$ en
conforme contre $1{,}4532$ en non conforme, soit un rapport de $0{,}9984$ alors que les niveaux
montent de $46\,\%$. L'invariance d'échelle du facteur entre états est donc établie à
$0{,}16\,\%$ près, et pour toute amplitude. Mais la dérive mesurée \emph{n'est pas} une pure
variation d'échelle : la modéliser réattribue à l'échelle ce que l'ajustement stationnaire lisait
comme de la forme, et l'indice de queue tombe de $0{,}5979$ à $0{,}5273$. Or c'est la forme qui
pèse. La composante de forme déplace seule le facteur de $-6{,}1\,\%$, soit $4{,}09$ écarts-types
de simulation, quand celle d'échelle ne le déplace que de $-0{,}2\,\%$, non résolu ; les deux
sont additives, le résidu d'interaction valant $0{,}23$ écart-type. Le facteur entre états passe
ainsi de $3{,}344$ à $3{,}124$, un déplacement résolu.

\paragraph{L'écart en euros, lui, bouge à peine, mais par compensation et non par
insensibilité.} Il passe de $13\,801$ à $13\,542$~M\euro, soit $-1{,}9\,\%$ pour un bruit de
simulation de $731$~M\euro, donc un déplacement non résolu. Il serait pourtant faux d'en conclure
que l'écart est robuste à la sévérité : la composante de forme seule le ramènerait à
$9\,393$~M\euro{} et celle d'échelle seule le porterait à $20\,041$. L'écart est donc très
sensible à la sévérité prise \emph{composante par composante}, et presque insensible à la seule
combinaison que la dérive produit. Ce qui survit sans réserve est le signe et
l'ordre, avec une marge large, et la direction est favorable : sous la sévérité dérivée
l'écart comme le facteur sont plus \emph{petits}, donc le chiffre publié n'est pas gonflé par la
dérive. La structure est celle du taux de dépassement gelé, un même défaut rendant le niveau
anti-conservateur et la thèse légèrement prudente, et seul le premier engage la solvabilité.

\paragraph{Le test se poursuit sur la charge annuelle agrégée, et il ne donne pas le même
verdict.} Tout ce qui précède porte sur les lois marginales prises séparément. Or le capital
n'est ni l'une ni l'autre : c'est le quantile de la charge annuelle. Un modèle composé peut avoir
deux marginales correctes et un agrégat faux, et c'est le reproche que cette classe de modèles
reçoit le plus souvent. L'observable noté est la somme annuelle des sinistres dépassant le seuil
d'apprentissage, qui est l'objet que la chaîne publiée produit puisqu'elle donne une sévérité
nulle sous le seuil ; la prédictive est obtenue par $200\,000$ années simulées à chaque origine.
Le verdict est négatif. La transformée intégrale de l'agrégat vaut $0{,}734$ en moyenne pour
$0{,}500$ attendu et Kolmogorov-Smirnov la rejette ($D = 0{,}484$, $p = 0{,}0042$) ; l'intervalle
de prévision annoncé à $90\,\%$ n'en couvre que $75{,}0\,\%$, avec trois dépassements de sa borne
haute concentrés dans les quatre dernières années notées.

\paragraph{Le rejet n'est pourtant pas structurel : c'est la dérive, vue sur l'objet qui porte le
capital.} Une dérive laisse une signature reconnaissable, puisqu'elle ne désaligne pas les
premières années notées mais seulement les dernières, l'apprentissage prenant du retard à mesure
qu'elle s'accumule. C'est ce que la mesure donne : la transformée intégrale vaut $0{,}611$ sur les
six premières années notées contre $0{,}857$ sur les six dernières, un écart que le test de
Mann-Whitney situe à $p = 0{,}0130$. La charge observée passe en moyenne de $363$ à
$1\,255$~M\euro{} entre les deux moitiés, quand la médiane prédictive ne passe que de $183$ à
$248$. Aucune limite nouvelle ne s'ajoute donc à l'inventaire du
chapitre~\ref{ch:robustesse} : la même s'y lit plus fort, l'agrégat cumulant la dérive de
l'échelle et celle du taux de dépassement quand chaque marginale n'en portait qu'une.

\paragraph{Un résultat inattendu au passage : la loi de comptage est décisive sur les comptes et
presque immatérielle sur la charge.} La binomiale négative gagne nettement sur les comptes. Notée
au score de probabilité continu classé sur la charge agrégée, son avance sur le Poisson tombe à
$0{,}3\,\%$. Le motif est celui que le mémoire avance déjà pour son propre quantile : la charge
annuelle est portée par un \emph{sinistre unique}, donc par la queue de sévérité et non par la
surdispersion du comptage. La binomiale négative reste le bon choix, mais elle cesse d'être un
déterminant du niveau de capital et il ne faut pas la présenter ainsi.

\paragraph{L'indépendance entre fréquence et sévérité cesse d'être une hypothèse non testée, et
elle tient.} Tout modèle composé la suppose, et une dépendance positive rendrait la queue de
l'agrégat plus lourde que celle que le modèle produit. Sur l'ensemble de la distribution, la
régression du logarithme du montant d'un sinistre sur le nombre d'incidents de son année donne
une pente négative et significative, $-0{,}02397$ à $p = 0{,}0419$, le millésime étant tenu
constant pour ne pas confondre l'effet avec la dérive. Mais sur les seuls \emph{excès}, qui sont
la population que le modèle tire, cette pente tombe à $-0{,}00331$ et cesse d'être significative
($p = 0{,}84$), soit $0{,}14$ de sa valeur. L'effet vit donc dans le bas de la distribution, où il
s'explique par la profondeur de collecte, une année à beaucoup d'incidents enregistrés étant une
année où l'on descend plus bas, et non par une propriété du risque. Prise au pied de la lettre
sur les excès, sa conséquence sur le quantile de l'agrégat est de l'ordre du pour cent et
symétrique en signe. Les corrélations de rang entre le compte annuel et la médiane ou le maximum
de l'année ne rejettent rien non plus, et le rang est ici le bon outil : à $\hat\xi$ voisin de
$0{,}6$ la variance de la sévérité est infinie, si bien qu'un test adossé à la variance de
l'agrégat n'aurait pas d'objet, du même argument d'existence que celui qui écarte la mesure de
queue sous la calibration PRC.

% =====================================================================
\section{Le corpus de post-mortems : les deux figures}
\label{sec:ann-corpus}

% SOURCES-SCRIPTS: 53 59

Complément du \S\ref{sec:postmortems}.

\begin{figure}[htbp]
  \centering
  \figover{Z17_postmortem_direction.png}
  \caption{La direction corroborée par codage de post-mortems
officiels, seule source où elle franchit le placebo. (a)~les $17$ transitions codées
$k\to j$ sur les sept rapports d'enquête publique du corpus initial, aucune n'allant à contre-sens
de la tendance dominante de sa paire ; le corpus étendu du \S\ref{sec:corpus-etendu} les porte à
vingt-deux sur dix rapports. (b)~le placebo directionnel est franchi ($z=+3{,}93$,
$p=0{,}0004$), là où le pas annuel d'OpRisk donnait $z=-0{,}33$ : c'est la qualité du codage,
non le volume, qui donne la puissance. (c)~P1 (gouvernance) est source pure ($9$
sorties, $0$ entrée) et P2 (gestion des incidents) est un puits ($0$ sortie, $7$ entrées) : la
structure que le classeur posait par jugement, retrouvée par un chemin indépendant. Le
confondant de narration reste déclaré et non levé.}
  \label{fig:postmortem}
\end{figure}

\begin{figure}[htbp]
  \centering
  \figover{Z20_corpus_etendu_pij.png}
  \caption{Le corpus étendu et la matrice ordonnée $p_{jk}$.
\textbf{(a)}~la matrice complète des médianes a posteriori, ligne émettrice et colonne
réceptrice : la ligne de $P_2$ et la colonne de $P_1$ sont vides, ce qui est la traduction
chiffrée du puits et de la source pure. (b)~les dix premiers liens ordonnés avec leur
crédible à $90\,\%$ : la direction ressort, l'amplitude reste large et aucun crédible ne passe
entièrement au-dessus de $0{,}5$. (c)~sensibilité à la sous-documentation des
défaillances non propagées : le niveau s'effondre, la part de direction tient.}
  \label{fig:corpus-etendu}
\end{figure}

% =====================================================================
\section{L'échelle de repli sur la direction : la reparamétrisation}
\label{sec:ann-echelle-u}

% SOURCES-SCRIPTS: 34

Complément du \S\ref{sec:echelle-u}.

La décomposition $W = S + A$ borne la direction en bloc. On peut faire mieux, en suivant
la méthode d'attaque puis de repli : reparamétrer la propagation par
\[
  \operatorname{logit} p_{ij} \;=\; \operatorname{logit} p_j + u_{ij},
\]
où $p_j$ est l'attractivité de la cible (ce qui subsiste quand la source ne compte
pas) et $u_{ij}$ l'excès propre à la source, c'est-à-dire la direction. La
décomposition n'est identifiée qu'une fois posée la normalisation $\sum_{i\neq j} u_{ij}=0$
pour chaque $j$ : c'est une convention de repérage, et non une information, ce que l'on
vérifie numériquement (une translation compensée de $p_j$ et de $u$ laisse le SCR
strictement inchangé). On descend alors une échelle de crans, du plus ambitieux au repli :
$u$ libre, puis $u$ de rang $1$ ($u_{ij}=b_j(a_i-\bar a)$, qui capte $74\,\%$ de
la direction posée et resserre les bornes de $25\,\%$ au prix d'une hypothèse de structure),
puis $u$ élicité, puis $u=0$.


\begin{figure}[htbp]
\centering
\figover{Z5_echelle_repli.png}
\caption{L'échelle de repli sur $u_{ij}$. (a)~Les bornes de capital se resserrent quand on
recule d'un cran (hypothèse de structure). (b)~Un rang $1$ capte $74\,\%$ de la direction
posée. (c)~Ce que la direction change à la décision : au cran nul, les gains de remédiation
de P1 et P4 se rejoignent.}
\label{fig:echelle}
\end{figure}


\begin{figure}[htbp]
\centering
\figover{Z2_comparaison_regimes.png}
\caption{Poser $W$ ou le borner, comparés avant élicitation. (a)~Un point contre un socle
et des bornes. (b)~Seule la part supérieure dépend de l'élicitation. (c)~Le régime~A suit
le biais du panel, la bande du régime~B ne bouge pas. (d)~Ce que l'élicitation doit
départager, et rien de plus.}
\label{fig:regimes}
\end{figure}

% =====================================================================
\section{La sensibilité aux probabilités de propagation : le détail}
\label{sec:ann-sensibilite-propagation}

% SOURCES-SCRIPTS: 15 76 30 66 68 81 63

Complément du \S\ref{sec:sensibilite-propagation}.

\paragraph{Pourquoi cette sensibilité n'est pas conduite par indices de Sobol.} L'outil
standard de l'analyse de sensibilité globale décompose la \emph{variance} de la sortie et
attribue à chaque entrée la part qu'elle en explique \citep{Owen2014, IoossPrieur2019}. Cet
outil n'est pas disponible ici, et le motif n'est pas une question de moyens : la
variance de la charge annuelle n'existe pas. L'indice de queue de la sévérité publié au
chapitre~\ref{chap:socle} excède un demi, seuil au-delà duquel le moment d'ordre deux d'une
loi de Pareto généralisée diverge, et une somme aléatoire de telles sévérités hérite de cette
divergence. La variance empirique d'une simulation reste finie, mais elle ne converge vers
rien : elle croît avec le nombre d'années tirées. Un indice défini comme un rapport de deux
quantités infinies n'a donc pas de valeur à estimer. C'est le même argument d'\emph{existence}
que celui qui écarte la TVaR comme mesure de couverture au \S\ref{sec:allocation}, sous la
calibration PRC : un ordre de moment plus haut, et la même conclusion.

La sensibilité de ce mémoire est par conséquent conduite par des grandeurs qui ne réclament
aucun moment d'ordre deux : élasticités du tornado renormalisé, ablations en échelle, et
décomposition de Möbius d'un écart de quantiles. Cela ne concerne pas la valeur de
Shapley du \S\ref{sec:allocation}, dont la fonction caractéristique est un écart de capital et
non une variance conditionnelle : elle est définie sans hypothèse de moment.

\paragraph{Le niveau de propagation est le levier le \emph{moins} sensible des quatre.} Le
tornado déplace chaque entrée du modèle vers le bas et vers le haut, à structure fixée. Le gain
$g$ y va de $0{,}5$ à $1{,}0$, ce qui déplace le surcoût de $5\,517$ à $6\,579$~M\euro. Le
\emph{seuil} de la loi de queue, lui, le déplace de $4\,862$ à $12\,944$~M\euro{} entre les
percentiles $90$ et $80$, et la fréquence de $4\,056$ à $8\,684$~M\euro{} en la divisant puis en
la multipliant par $1{,}5$ ; la surdispersion, de $6\,449$ à $6\,834$. La fragilité du
capital n'est donc pas dans la propagation, elle est dans l'ajustement de valeurs extrêmes,
et c'est l'inverse de ce que l'on redoute spontanément d'un modèle de contagion. Sur les huit
variantes, le surcoût reste compris entre $4\,056$ et $12\,944$~M\euro, toujours positif et
toujours pluri-milliard.

\begin{attention}
\textbf{Ces niveaux sont ceux de la configuration de base du tornado} ($6\,640$~M\euro), qui
n'est pas l'écart à quatre canaux de $14\,139$~M\euro{} du \S\ref{sec:kpi} : ce tornado fait
varier un paramètre à la fois autour d'un état non conforme unique. Ce qui se lit ici est
l'\emph{amplitude relative des leviers}, en aucun cas un niveau de capital.
\end{attention}

\paragraph{Deux défauts de ce tornado, et une renormalisation qui les corrige.} Tel qu'il vient
d'être lu, ce tornado prête à deux objections, et toutes deux portent.

La première est de comparabilité. Les quatre leviers sont balayés sur des plages qui
n'ont pas la même vraisemblance : le gain va de $0{,}50$ à $1{,}00$, amplitude \emph{posée}, la
fréquence de la moitié au triple. Un tornado compare des longueurs de barre, or une barre est
longue parce que le paramètre bouge beaucoup ou parce qu'on l'a fait bouger beaucoup.

La seconde est plus sérieuse : le levier \og seuil\fg{} n'est pas un levier. Déplacer le
seuil réajuste la loi de queue, si bien que $\xi$, $\sigma$ et le taux de dépassement changent
avec lui : il déplace \emph{quatre} paramètres quand le gain en déplace un. En le décomposant, le
seuil pris seul ne déplace que $+320$~M\euro, l'indice de queue seul en déplace $-13\,993$, et
les composantes ne s'additionnent pas mais se compensent, leur somme signée valant
$-11\,517$ pour un paquet à $-8\,081$. La barre publiée sous le nom de \og seuil\fg{} est donc plus
courte que l'effet du seul indice de queue, tout en lui étant attribuée.

\paragraph{Ce qui survit à la renormalisation, et c'est l'essentiel.} La normalisation qui rend
les leviers comparables est l'élasticité, sans dimension et sans plage : on applique à
chacun la même perturbation relative de $10\,\%$, un paramètre à la fois et sans réajustement.

\begin{center}\footnotesize
\renewcommand{\arraystretch}{1.15}
\begin{tabular}{@{}l l r@{}}
\toprule
\textbf{Levier} & \textbf{Statut} & \textbf{Élasticité} \\
\midrule
Indice de queue $\xi$ & estimé & $3{,}79$ \\
Échelle $\sigma$ & estimé & $0{,}97$ \\
Taux de dépassement $p_u$ & estimé & $0{,}82$ \\
Fréquence $\lambda$ & estimé & $0{,}45$ \\
Sur-dispersion $\varphi$ & posé & $-0{,}44$ \\
\textbf{Gain de propagation} $\boldsymbol{g}$ & \textbf{posé} & $\mathbf{0{,}37}$ \\
Seuil $u$ seul & choix de méthode & $0{,}03$ \\
\bottomrule
\end{tabular}
\end{center}

\vspace{4pt}
\noindent \textbf{La conclusion tient.} La propagation
a l'élasticité la plus faible des leviers substantiels, $0{,}37$ contre $3{,}79$ pour l'indice de
queue, soit un rapport de dix. Ce classement ne dépend d'aucune convention d'affichage, alors
qu'une comparaison d'amplitudes pouvait être retournée en changeant une plage.
Deux précisions l'accompagnent. Le seuil \emph{pris seul} a une élasticité plus faible encore,
$0{,}03$, ce qui confirme que tout ce qu'on lui attribuait venait de l'indice de queue qu'il
entraîne. Et une élasticité est \emph{locale} : elle ne dit pas ce qui se passe loin du point de
base, où la réponse est l'invariance de l'ordre en $g$ du paragraphe suivant.

\paragraph{Une dernière asymétrie, qui n'est pas de calcul mais de statut.} Deux des leviers,
$g$ et $\varphi$, sont posés : ils n'ont aucun intervalle d'estimation, et la plage
qu'on leur donne est un choix d'affichage. Les autres sont estimés et ont une plage défendable :
sur son intervalle à $90\,\%$ publié, l'indice de queue déplace le surcoût de $1\,908$ à
$32\,666$~M\euro. Comparer la barre d'un paramètre posé à celle d'un paramètre estimé revient
donc à comparer une convention à une incertitude, et c'est une raison de plus de lire ce tornado
en élasticités.

\paragraph{Et la thèse ne dépend pas des trois valeurs posées.} Le capital est croissant en $g$,
et ce n'est pas supposé mais vérifié à nombres aléatoires communs : toute
% SOURCES-SCRIPTS: 66
correspondance entre états de conformité et valeurs de $g$ qui respecte l'ordre produit
l'écart, dès lors que le forfait réglementaire reste plat. L'élasticité d'arc du capital au gain
vaut $0{,}410$. Seule l'\emph{amplitude} de l'écart est un scénario ; son signe et son ordre ne
le sont pas.

\begin{attention}
\textbf{Ces trois valeurs ne sont pas ancrées sur une source de supervision, et c'est une
décision plutôt qu'une lacune.} La question se pose naturellement : les référentiels de maturité
publiés par les autorités nationales et européennes ne permettraient-ils pas de caler $g$ au lieu
de le poser ? La réponse est non, et le motif n'est pas la disponibilité mais la \emph{nature} de
ce qu'ils publient. Ils donnent des attentes prudentielles et des échelles de maturité, jamais des
probabilités de propagation entre domaines de contrôle. Passer des unes aux autres demanderait une
correspondance maturité vers gain de propagation, c'est-à-dire une seconde correspondance
posée, et le résultat serait strictement pire que la valeur posée d'aujourd'hui : il aurait
l'apparence d'un calibrage sans en être un.

\textbf{Et l'invariance rend l'ancrage sans objet.} Puisque le capital est croissant en $g$ et que
toute correspondance respectant l'ordre produit l'écart, un ancrage n'achèterait que
l'\emph{amplitude}, soit exactement la part que ce mémoire déclare comme un scénario. C'est le
même arbitrage que celui qui a écarté le recours au jugement d'expert formalisé : une source
faiblement calibrée introduirait une incertitude non déclarable en échange d'une ignorance
mesurée. Une valeur posée et annoncée comme telle vaut mieux qu'une valeur dérivée d'une échelle
qui ne mesure pas la grandeur voulue.
\end{attention}

\paragraph{La direction, elle, ne se stresse pas : elle se borne, et la borne se gradue.} Puisque
les sources disponibles renseignent la partie symétrique de $W$ et non sa partie antisymétrique
(proposition~\ref{prop:nonid}), la sensibilité à
la direction ne s'écrit pas comme un déplacement autour d'un point mais comme une largeur
d'ensemble, indexée par le degré d'ignorance $t$ que l'on assume sur cette direction. Les
$2^{10}$ sommets de l'ensemble admissible sont énumérés exhaustivement, et tous sont admissibles.

\begin{center}\footnotesize
\renewcommand{\arraystretch}{1.25}
\begin{tabular}{@{}l l r r@{}}
\toprule
\textbf{Ignorance $t$} & \textbf{Capital borné} & \textbf{Largeur} & \textbf{En part du socle} \\
\midrule
$t=0{,}25$ & $[7\,837\,;\,8\,174]$ & $336$ & $6{,}4\,\%$ \\
$t=0{,}50$ & $[7\,605\,;\,8\,300]$ & $695$ & $13{,}2\,\%$ \\
$t=0{,}75$ & $[7\,348\,;\,8\,431]$ & $1\,083$ & $20{,}5\,\%$ \\
$t=1{,}00$ (état actuel) & $[\mathbf{6\,858}\,;\,\mathbf{8\,697}]$ & $\mathbf{1\,839}$ & $\mathbf{34{,}9}\,\%$ \\
\bottomrule
\end{tabular}
\end{center}

\vspace{4pt}
\noindent En M\euro. La lecture est la suivante, et elle est plus favorable qu'un intervalle de
confiance ordinaire : l'ignorance sur la direction se paye linéairement, et son coût
maximal est connu d'avance. À ignorance totale le capital reste enfermé dans un facteur
$1{,}27$ entre ses bornes, alors que le socle sans aucune contagion vaut $5\,275$~M\euro{} : la
contagion ajoute entre $30$ et $65\,\%$ à ce socle, et l'on sait exactement de combien cette
fourchette rétrécirait si la direction était mieux connue. À titre de repère, la direction nulle
($W$ symétrique) donne $7\,993$~M\euro{} et le point du classeur d'expert $8\,110$, soit
respectivement $+51{,}5\,\%$ et $+53{,}7\,\%$ au-dessus du socle : le point d'expert
est un point de cet ensemble, pas sa mesure, et l'essentiel du surcoût de contagion ne vient pas
de la direction mais de la co-occurrence, qui est identifiée.

\paragraph{Enfin, la décomposition de l'écart a son propre bruit, et il est publié.} Les quinze
termes qui répartissent l'écart de capital entre les quatre canaux sont accompagnés de leur
écart-type de simulation, et neuf sur quinze seulement ont un signe résolu sur seize
graines (\S\ref{sec:kpi}). Ce n'est pas un défaut de la mesure, c'est sa précision : une
interaction est une différence de différences, l'estimateur le plus fragile du dispositif, et
les termes d'ordre trois et quatre figurent dans la table parce que la réconciliation l'exige,
non parce qu'ils sont mesurés.

\paragraph{La même conclusion par un chemin indépendant : l'ablation en échelle.} Un tornado
\emph{perturbe} un paramètre ; une ablation \emph{retire} une brique du mécanisme. Les deux
protocoles n'ont en commun ni la question posée ni la métrique, de sorte qu'ils peuvent
diverger. Ils ne divergent pas. Retirer la propagation, en posant $g=0$,
% SOURCES-SCRIPTS: 81
fait tomber la VaR de $20\,188$ à $16\,086$~M\euro, soit $-20\,\%$ ; retirer la queue lourde en
la ramenant à une loi légère la fait tomber de $76\,\%$. La brique la plus lourde du
modèle est la queue de sévérité, pas la contagion, et le facteur entre les deux vaut $3{,}7$.
C'est mot pour mot la conclusion du tornado renormalisé, obtenue sans le renormaliser.

Deux précisions que ce protocole impose, et la seconde n'était pas attendue. Le contrôle est
exact : le moteur ventilé par pilier reproduit l'état de référence à $0{,}0000$~M\euro{} près,
donc les colonnes se somment bien à la charge agrégée. Et la forme de queue est héritée,
non produite : le rapport $\mathrm{CTE}/\mathrm{VaR}$ vaut $2{,}075$ à la référence et ne bouge
pratiquement pas quand on retire la propagation ($2{,}061$). La cascade déplace donc le
\emph{niveau} du capital sans façonner la \emph{forme} de sa queue, qu'elle reçoit de la
sévérité. Une lecture qui attribuerait la lourdeur de queue au mécanisme de contagion serait
fausse, et c'est une lecture naturelle.

\begin{attention}
\textbf{Deux ordres de piliers sortent de cette ablation, et un seul est un signal.} L'ordre par
pilier \emph{d'amorce}, $P4 > P1 > P2 > P3 > P5$, est stable entre graines. L'ordre par pilier
\emph{touché} ne l'est pas, et il ne peut pas l'être : les parts touchées tiennent toutes entre
$17$ et $23\,\%$, écart bien inférieur au bruit d'un quantile porté par une sévérité de variance
infinie. Le second est reporté pour mémoire et n'autorise aucune conclusion. C'est la même
limite que celle déjà déclarée sur l'allocation d'Euler au \S\ref{sec:allocation}, retrouvée ici
par un autre chemin.
\end{attention}

\paragraph{Et sur ce classement, le préprint de cascade climatique dit l'inverse : il faut le
dire, et l'expliquer.} Le préprint cité au chapitre~\ref{ch:positionnement} est invoqué trois fois
dans ce mémoire comme convergent, sur la séparation des échelles, le protocole à marges appariées
et l'expérience d'ablation. Sur le \emph{classement des leviers}, il conclut à l'opposé, et taire
ce point serait retenir d'une source ce qui arrange.

\begin{center}\footnotesize
\renewcommand{\arraystretch}{1.2}
\begin{tabular}{@{}l l l@{}}
\toprule
& \textbf{leur cascade} & \textbf{celle-ci} \\
\midrule
Brique la plus lourde à l'ablation & propagation dirigée, $-53{,}9\,\%$ & queue, $-76\,\%$ \\
Propagation & dominante & $-20\,\%$ \\
Tête du tornado & probabilité d'arête, $-25{,}0$ à $+26{,}6\,\%$ & indice de queue, élasticité
$3{,}79$ \\
Propagation au tornado & première barre & élasticité $0{,}37$, dernière \\
\bottomrule
\end{tabular}
\end{center}

\vspace{4pt}
\noindent Leur conclusion nomme la propagation dirigée et la fréquence annuelle comme les
principaux moteurs ; la présente section nomme l'ajustement de valeurs extrêmes. Les chiffres de
leur travail sont enregistrés sous le statut de \emph{citation externe, non recalculable}, au
même titre que les autres sources non versionnées.
% SOURCES-SCRIPTS: 63

\begin{cle}
\textbf{Ce n'est pas un désaccord de mesure, et la cause est identifiable dans leur propre
texte.} Leur sévérité est bornée : réponse de sévérité bornée, perte plafonnée par une
transformation à capacité, multiplicateurs lognormaux de moyenne un et d'écart-type logarithmique
$0{,}10$ et $0{,}15$. Elle \emph{n'a pas d'indice de queue}, et leur ablation ne contient donc
aucune brique de queue : on ne retire pas ce qui n'est pas là. Le paramètre le plus proche qu'ils
font varier est la raideur de la réponse en sévérité, et elle sort derrière la propagation. Ici,
la sévérité est une GPD de variance infinie et la queue domine tout.

\textbf{Les deux classements sont donc corrects, chacun dans son modèle, et ce qui les sépare est
l'indice de queue et non l'architecture.} La conséquence de citation est la même que celle déjà
posée pour l'échelle des quantiles au \S\ref{sec:benchmark} : ce préprint se cite sur le
\emph{protocole}, jamais sur l'ordre d'un résultat en queue. Un lecteur qui ouvrirait
le préprint et y trouverait la propagation en tête n'a donc pas trouvé une contradiction, il a
trouvé un modèle sans queue lourde.
\end{cle}

% =====================================================================
\section{Ce que la donnée ouverte permet de piloter}
\label{sec:ann-regard-reglementaire}

% SOURCES-SCRIPTS: 28 27

Complément du \S\ref{sec:regard-reglementaire}.

\begin{cle}
\textbf{Jusqu'où l'open data descend, et où il s'arrête.} Un test sur
% SOURCES-SCRIPTS: 28
la base VCDB montre que les sous-piliers de P2 (familles d'action) et, plus grossièrement,
la détection de P3 sont exploitables sur donnée ouverte (plusieurs centaines d'incidents en
finance, $45\,\%$ de détection lente). Mais le pilier des tiers P4 s'effondre : $56$
incidents à acteur partenaire en finance, et la donnée ouverte sous-code notoirement
l'externalisation. Le pilotage fin de P4 exige donc une donnée que seul un registre
d'externalisation, tenu par l'entité ou le régulateur, peut fournir. C'est le même verdict
que celui du chapitre~\ref{ch:identifiabilite}, décliné au niveau des sous-piliers.
\end{cle}

\begin{figure}[htbp]\centering
\figover{W_benchmark_sf.png}
\caption{Ce que chaque cadre exprime de l'effet DORA : $0\,\%$ pour la Formule Standard,
$62\,\%$ pour la copule, $100\,\%$ pour la cascade.}
% SOURCES-SCRIPTS: 27
\label{fig:benchmark-sf}
\end{figure}

% =====================================================================
\section{Un ordre de grandeur à l'épreuve du réel : le détail}
\label{sec:ann-ancrage}

% SOURCES-SCRIPTS: 60 49

Complément du \S\ref{sec:ancrage}.

\paragraph{Contre la Formule Standard, sur la même entité.} Pour l'entité notionnelle
($15\,000$~M\euro{} de provisions), la branche provisions de la charge opérationnelle vaut
$0{,}03\times 15\,000 \approx 450$~M\euro{} si ces provisions sont non-vie. Ce repère n'est pas la
charge réglementaire, qui retient le maximum des branches primes et provisions sous le plafond
assis sur le BSCR, et il n'en est pas une borne ; il en donne l'ordre de grandeur. Le socle de la cascade
($5\,275$~M\euro, chapitre~\ref{chap:partielle}) lui est une dizaine de fois
supérieur. Deux lectures cohabitent, et aucune ne se tranche sans donnée d'entité : soit la
charge générique \emph{sous-capte massivement} le cyber, ce que le benchmark
ci-dessus établit par ailleurs ; soit la sévérité OpRisk, tirée de grandes institutions
américaines, \emph{sur-représente} pour une entité de taille moyenne. Les deux jouent, à des
degrés qu'on ne sait pas départager.

\paragraph{Contre le marché.} Un capital cyber de l'ordre de $5$~milliards d'euros pour une
\emph{seule} entité représenterait une fraction substantielle des primes cyber annuelles
mondiales, elles-mêmes de l'ordre de la dizaine de milliards d'euros
\citep{munichre_cyber_2024}. Le rapport est invraisemblable à l'échelle d'un assureur moyen :
il signale que le niveau, transporté d'une sévérité de grandes institutions vers une entité
notionnelle, ne saurait être lu comme le capital d'une entité réelle.

\paragraph{Au niveau du sinistre unique, l'ordre de grandeur est en revanche crédible.} Le
constat précédent porte sur l'\emph{agrégat} annuel. À l'échelle du \emph{sinistre unique},
les quantiles du modèle tombent dans la fourchette des grandes pertes cyber réelles
(figure~\ref{fig:benchmark-externe}) : le quantile de sévérité $\qsev_{99{,}5\,\%}(X)$
($663$~M\euro) et sa moyenne de queue $\qbarsev_{99{,}5\,\%}(X)$ ($1\,752$~M\euro) encadrent
Merck-NotPetya ($\sim$1,3~Md\euro) et Equifax ($\sim$1,3), et restent sous Change
Healthcare 2024 ($\sim$2,2) comme sous le choc systémique NotPetya ($\sim$9~Md\euro
au total). La brique de sévérité est donc crédible \emph{à l'échelle de l'événement} ; c'est
son agrégation à une fréquence systémique ($\sim$21,6 incidents matériels par an, à
l'échelle du secteur financier mondial) qui porte l'agrégat annuel, de l'ordre du marché
cyber mondial ($\sim$15~Md\euro de primes en 2025), à une échelle de \emph{secteur}, non
d'entité isolée. Le modèle n'est donc ni gonflé (il colle aux pertes réelles au niveau
unitaire) ni aveugle (il capte, au niveau agrégé, le risque que la Formule Standard manque).

\paragraph{La descente d'échelle rend l'argument vérifiable.} Le diagnostic qui précède se
teste : si l'écart tient bien à la fréquence d'entrée et non à la mécanique de cascade, alors
il suffit de \emph{redescendre} $\lambda$ au niveau d'une entité, à matrice $W$ et moteur de
propagation inchangés, pour retrouver un ordre de grandeur d'entité. Deux canaux, et non un
seul, séparent les deux échelles, et il faut les composer.

\paragraph{Canal 1, la fréquence, lue à la taille de la cible.} Le panel firme-année du SAS
OpRisk fournit deux lectures immédiates, $\lambda = 0{,}10$ sur toutes les firmes et
$\lambda = 0{,}21$ sur celles à au moins dix événements. Aucune des deux n'est la bonne, pour
deux raisons. La première tient à la taille : les firmes du second groupe ont
$372$~Md\$ d'actifs médians, soit dix-neuf fois l'entité notionnelle, et leur taux ne
saurait donc la décrire. La seconde tient à la sélection : définir un groupe par
\og{}au moins dix événements\fg{} puis calculer son nombre moyen d'événements, c'est
conditionner sur la grandeur même que l'on estime, ce qui la biaise nécessairement vers le
haut. On estime donc $\lambda$ comme une \emph{fonction} de la taille, par une binomiale
négative à lien logarithmique sur le panel, et on la lit à la taille de la cible :
l'élasticité vaut $b_\lambda = 0{,}074$ (IC~$95\,\%$ $[0{,}055\,;0{,}094]$, $z = 7{,}6$), d'où
$\lambda_{\text{cible}} = 0{,}092$.

\paragraph{Canal 2, la sévérité, à la même taille.} L'élasticité sévérité/taille du
\S\ref{sec:biais-chiffres} ($b = 0{,}087$, estimée par vraisemblance tronquée) transpose de la
firme médiane de la base ($122$~Md\$ pondérés par événement) à la cible et donne un
multiplicateur de $0{,}854$. La sévérité GPD étant une famille d'échelle au-dessus du seuil et
la source OpRisk n'étant pas plafonnée, cette homothétie déplace le quantile exactement dans
le même rapport.

Trois enseignements, dont un qui corrige une lecture trop rapide.

D'abord, la non-linéarité. Diviser $\lambda$ par $235$ divise la perte moyenne par $270$,
c'est-à-dire presque à proportion, mais le capital par $48$ seulement. C'est la signature du
\emph{principe du grand saut unique} : la moyenne est portée par le nombre d'événements, le
quantile extrême par l'ampleur du plus grand. À cette fréquence, la quasi-totalité des années
sont vides et le quantile $99{,}5\,\%$ ne décrit plus une accumulation mais un sinistre isolé
suivi de sa cascade. Le capital d'entité n'est donc pas le capital de secteur divisé
par le nombre d'entités, et le lire ainsi conduirait à le sous-estimer nettement.

Ensuite, la coïncidence avec la Formule Standard
ne survit pas à la correction. Lue au seuil \og{}$\geq 10$ événements\fg{}, l'entité ressortait
à $488$~M\euro, soit l'ordre de grandeur exact du repère forfaitaire ($450$~M\euro), et
l'on pouvait croire l'écart refermé. Une fois la fréquence lue à la bonne taille et la
sévérité transposée, le modèle donne $169$~M\euro, soit $0{,}38$ fois ce repère. La
convergence apparente venait donc du \emph{choix du seuil}, pas du modèle. La conclusion à en
tirer n'est pas que le modèle est faux : c'est que le niveau de ce repère ne dit rien
et que seule son insensibilité est en cause, ce que le paragraphe suivant chiffre.

Enfin, la lecture d'entité doit être une bande. La direction de $W$ n'étant pas
identifiée (chapitre~\ref{ch:identifiabilite}), le $169$~M\euro{} est un point de l'ensemble
admissible, pas sa mesure. En énumérant les $1\,024$ sommets du pavé à $t=1$, le SCR d'entité
tient dans $[131{,}5\,;183{,}3]$~M\euro{}, au-dessus d'un socle sans contagion de
$78{,}1$~M\euro{} qui, lui, ne dépend d'aucune hypothèse directionnelle et vaut à lui seul
$43\,\%$ de la borne haute. Au sens de la part \emph{non exposée} à l'ignorance
directionnelle employé au chapitre~\ref{chap:partielle}, c'est-à-dire un moins le rapport de
la largeur de bande à la borne haute, l'échelle d'entité fixe $72\,\%$ du capital, contre
$79\,\%$ à l'échelle du secteur : la descente d'échelle élargit la part exposée, sans la
faire basculer.

Une dernière précaution, et elle est lourde : ce $169$~M\euro{} est une borne supérieure
d'ordre de grandeur, non une charge plausible. Le \S\ref{sec:entites-reelles} établit sur des
bilans réels que la transposition a une borne inférieure de taille, et que l'entité notionnelle
se situe en dessous. Rien de ce qui suit dans ce chapitre ne repose sur le niveau lui-même, mais
sur les \emph{écarts} et les \emph{rapports} qui s'en déduisent ; il faut néanmoins l'avoir lu
avant de citer le chiffre.

\begin{figure}[htbp]\centering
\figover{S21_descente_echelle.png}
\caption{La descente d'échelle, rendue cohérente. (a)~$\lambda$ estimé comme une
fonction de la taille de firme sur le panel firme-année, et lu à la taille de l'entité visée
plutôt que choisi dans un seau ; le trait pointillé rappelle le seau $0{,}21$, qui correspond
à des firmes dix-neuf fois plus grandes. (b)~les deux corrections composées, et le
repère de Formule Standard, dont la coïncidence apparente avec l'ancien chiffre ne survit pas.
\textbf{(c)}~le résultat d'entité est une bande sur les $1\,024$ sommets admissibles, non un
point.}
% SOURCES-SCRIPTS: 60
\label{fig:descente}
\end{figure}

\begin{attention}
\textbf{Ce que la correction d'échelle déplace, et ce qu'elle laisse intact.} La descente à la
bonne taille divise les niveaux d'entité par près de trois. Il faut donc dire lesquelles des
conclusions en dépendent. Le multiple de capital ($52{,}0 \to 46{,}5$) et le coût de
détention rapporté à la sinistralité attendue ($2{,}47 \to 2{,}21$) bougent de moins de
$11\,\%$ : leurs deux termes se rapportent au même SCR et se déplacent ensemble, ce qui
vérifie sur un cas non préparé pour cela la thèse du \S\ref{sec:detenir-transferer}. En
revanche le gain du dimensionnement du traité passe de $-47\,\%$ à $-64\,\%$, et cela
doit être dit plutôt que rangé avec les rapports : à fréquence plus basse, la perte annuelle
est davantage dominée par un sinistre unique, une part plus grande de l'espérance se situe
au-delà du SCR et échappe à une portée bornée par lui, de sorte que la part cédée tombe de
$61$ à $38\,\%$. Le \emph{sens} de la conclusion est renforcé, son \emph{amplitude} chiffrée
dépend de l'échelle et ne doit pas être citée comme une grandeur d'entité.
\end{attention}

\paragraph{Pourquoi la mise à l'échelle n'est pas proportionnelle, et c'est mesuré.} La façon la
plus naturelle de descendre un capital sectoriel à une entité est de le multiplier par le poids
de cette entité dans le secteur. C'est le repère contre lequel la méthode se distingue, et il
faut dire pourquoi elle s'en écarte plutôt que le passer sous silence. Une renormalisation
proportionnelle affirme une chose précise : que le besoin de capital est d'élasticité
un à la taille. Les deux canaux par lesquels la taille agit ont été estimés séparément, et
aucun des deux ne s'en approche.

\begin{center}\footnotesize
\renewcommand{\arraystretch}{1.25}
\begin{tabular}{@{}l l l@{}}
\toprule
\textbf{Canal} & \textbf{Élasticité à la taille} & \textbf{Ce que dit la mesure} \\
\midrule
Fréquence & $+0{,}0744$ (ET $0{,}0098$, $z=+7{,}59$) & très significative, et très loin de $1$ \\
Sévérité & $0{,}087$, IC $[0{,}026\,;0{,}148]$ & contient presque zéro \\
\bottomrule
\end{tabular}
\end{center}

\vspace{4pt}
\noindent Chacun des deux est \emph{d'un ordre de grandeur} sous l'unité que la proportionnalité
suppose, et le second est le plus frappant : la sévérité ne se met pratiquement pas à l'échelle
de la taille, d'où le multiplicateur de $0{,}8545$, soit $-14{,}6\,\%$ seulement, appliqué en
descendant de la population de calibration à l'entité cible. La méthode retenue est donc
à deux canaux séparés, la fréquence par le lien logarithmique du modèle de comptage et
la sévérité par son élasticité mesurée, et non un facteur unique.

\begin{attention}
\textbf{Les deux approches se trompent en sens opposés.} Une descente proportionnelle \emph{sous}-estimerait le besoin d'une petite entité,
puisqu'elle rétrécirait la sévérité alors que la donnée dit qu'elle ne rétrécit presque pas. Le
modèle retenu fait l'inverse : en ne plafonnant pas la sévérité à l'échelle d'entité, il affirme
qu'une entité de quelques milliards subit à peu près la sévérité d'une institution mondiale, ce
qui \emph{majore}. C'est exactement pourquoi le chiffre d'entité est publié comme une
borne supérieure d'ordre de grandeur et pourquoi la méthode déclare une borne inférieure
de validité en taille (\S\ref{sec:entites-reelles}). Substituer la proportionnalité ne
supprimerait pas le problème, elle en changerait le signe sans le déclarer.
\end{attention}

\begin{figure}[htbp]\centering
\figover{J5_benchmark_externe.png}
\caption{Benchmark externe : la confrontation au réel, à deux échelles
qu'il ne faut pas confondre. (a)~au niveau du \emph{sinistre unique}, les quantiles du modèle
(quantile de sévérité $663$~M\euro{} et moyenne de queue $1\,752$~M\euro{} à $99{,}5\,\%$, et non
des capitaux) s'insèrent parmi les grandes pertes
cyber réelles, encadrant Merck-NotPetya, Equifax et Change Healthcare tout en restant sous le
choc systémique NotPetya : la brique de sévérité est crédible. (b)~à l'échelle
\emph{annuelle}, le capital est d'ordre systémique, comparable au marché cyber mondial, et
domine de loin la charge de Formule Standard. C'est l'agrégation à une fréquence sectorielle,
non la sévérité, qui porte l'agrégat à une échelle de secteur.}
% SOURCES-SCRIPTS: 49
\label{fig:benchmark-externe}
\end{figure}

% =====================================================================
\section{L'agrégation au capital total : le contrefactuel et le plafond}
\label{sec:ann-agregation}

% SOURCES-SCRIPTS: 95 65

Complément du \S\ref{sec:agregation-scr}.

\paragraph{Le contrefactuel, et ce qu'il mesure.} Pour dire ce que cette convention implique, on
agrège le même montant $D$ comme un module de plus, corrélé à $\rho$ avec le BSCR pris
comme un bloc. Le taux de reconnaissance vaut alors
\[
  \tau(\rho, d) \;=\; \frac{\sqrt{1 + 2\rho d + d^{2}} - 1}{d},
  \qquad d = \frac{D}{\mathrm{BSCR}},
\]
et son développement en $d = 0$ donne $\tau(\rho, d) = \rho + \tfrac{d}{2}\,(1 - \rho^{2}) +
O(d^{2})$. Un montant petit devant le BSCR est donc reconnu à hauteur de $\rho$,
et de $\rho$ seulement, quand la convention additive le reconnaît en entier. L'écart entre les
deux traitements vaut $1 - \rho$ : il ne dépend pas de la taille, et il ne se referme qu'à
corrélation parfaite. Ce contrefactuel n'est pas ce que fait Solvabilité~II ; il sert à
\emph{mesurer} la convention en vigueur, jamais à en proposer une autre.

\paragraph{Ce que la donnée publiée suffit à établir.} Le BSCR ne figure pas parmi les
champs retenus, et il n'a pas besoin d'y figurer : $\mathrm{Adj} \le 0$ donne
$\mathrm{BSCR} \ge \mathrm{SCR} - \mathrm{SCR}_{Op}$, et le plafond donne
$\mathrm{BSCR} \ge \mathrm{SCR}/1{,}3$. Minorer le BSCR \emph{majore} $\tau$ : les taux
ci-dessous sont donc des majorants, et l'écart entre les deux conventions vaut au moins ce qu'on
y lit.


\paragraph{Le plafond, qui ajoute une borne à la cécité déjà établie.} La charge forfaitaire est
plafonnée à $0{,}3 \times \mathrm{BSCR}$, quelle que soit l'exposition opérationnelle réelle. Sur
l'assureur non-vie~A, dont l'état \textsc{s.25.01} publie un module opérationnel de
$59{,}0$~M\euro, ce plafond vaut $109{,}8$~M\euro{} et la marge restante $50{,}8$~M\euro. Le
besoin calculé, $112{,}1$~M\euro, en représente $2{,}21$ fois, et l'écart entre états,
$34{,}9$~M\euro, en représente $0{,}69$ fois. Autrement dit, même un forfait rendu
sensible à la conformité saturerait avant d'avoir exprimé ce besoin : la Formule Standard n'est
pas seulement insensible, elle est bornée. La lecture $\mathrm{Adj} = 0$ est ici \emph{posée et
déclarée}, un ajustement strictement négatif relevant le BSCR, donc le plafond, donc la
marge ; l'état \textsc{s.25.01} de l'entité permet de la lever au cas par cas.

% =====================================================================
\section{Détenir ou transférer : le dimensionnement du traité}
\label{sec:ann-detenir}

% SOURCES-SCRIPTS: 58 67 60

Complément du \S\ref{sec:detenir-transferer}.

\paragraph{Le dimensionnement, plutôt que le tout ou rien.} La question n'est pas binaire : la
capacité du marché est bornée et l'on ne peut pas tout céder. On dimensionne donc un traité en
excès de perte annuel de portée $L$, en minimisant le coût total
\[
  \mathrm{co\hat ut}(L) \;=\; \underbrace{\mathbb E\bigl[\min(X,L)\bigr]\big/\,0{,}485}_{\text{prime}}
  \;+\; \mathrm{CoC}\times \underbrace{\mathrm{VaR}_{99{,}5\%}\bigl(X-\min(X,L)\bigr)}_{\text{capital résiduel}} .
\]
Le capital résiduel vaut $\max(\mathrm{SCR}-L,0)$, de sorte que l'optimum tombe exactement en
$L^\star=\mathrm{SCR}=169$~M\euro{} : au-delà, la portée ne rachète plus aucun capital. Le
coût total passe de $8{,}0$ à $2{,}9$~M\euro{} par an, soit $-64\,\%$, en
cédant $38\,\%$ de la sinistralité. Le mécanisme est transparent : près de l'optimum, un euro
de portée supplémentaire coûte $1{,}04\,\%$ de prime et libère $4{,}75\,\%$ de coût du capital.

\begin{attention}
\textbf{Ce gain est une borne supérieure, et il faut le citer comme telle.} Trois raisons
empêchent de conclure au chiffre brut, et toutes vont dans le même sens. La capacité :
une couverture agrégée de l'ordre de $170$~M\euro{} reste plaçable, mais elle mobilise déjà
plusieurs lignes sur un marché cyber dont la capacité par risque est bornée. Le risque
de base : les polices excluent la guerre et les
infrastructures essentielles, donc une part de la queue modélisée n'est pas assurable. La
\emph{contrepartie} : sous Solvabilité~II, céder ne supprime pas l'exigence, cela substitue
au risque cyber un risque de défaut de contrepartie sur les créances de réassurance. En
résumant les trois par un coefficient d'inefficacité $\kappa$, le capital résiduel devient
$\max(\mathrm{SCR}-L,0)+\kappa L$, et le transfert reste avantageux tant que le taux marginal
de prime reste sous $\mathrm{CoC}\times(1-\kappa)$. Le seuil critique vaut
$\kappa^\star=1-0{,}0103/0{,}0475=78\,\%$ : il faudrait que plus des trois quarts de
la couverture soient inefficaces pour que la détention redevienne préférable. La conclusion
qualitative est donc robuste, le chiffre ne l'est pas.
\end{attention}

\paragraph{Ce que la conformité change à l'arbitrage.} La remédiation agit sur la fréquence
d'amorce. En rejouant l'arbitrage à $\lambda$ réduit, ce qui simule une entité devenue
conforme, la portée optimale suit le capital ($169 \to 97 \to 43$~M\euro{} quand $\lambda$
passe de $0{,}092$ à $0{,}064$ puis $0{,}046$) et le coût total tombe de $2{,}9$ à
$0{,}5$~M\euro{} par an, soit $-81\,\%$. \textbf{La conformité ne fait
% SOURCES-SCRIPTS: 67
donc pas que réduire le
capital : elle déplace le point d'équilibre entre rétention et transfert}, et réduit
simultanément les deux termes du coût. C'est un résultat de \emph{pilotage} et non de mesure,
et il met les trois leviers de l'entité (remédier, détenir, transférer) sur un seul axe en
euros par an, ce qu'aucune des comparaisons de niveau ne permettait.

\begin{figure}[htbp]
  \centering
  \figover{S19_detenir_ou_transferer.png}
  \caption{Détenir ou transférer : le SCR confronté au prix de
marché du même risque, à l'échelle d'entité corrigée. (a)~le coût annuel du capital
($8{,}0$~M\euro{} au taux révisé de $4{,}75\,\%$) vaut $2{,}2$ fois la sinistralité attendue
($3{,}6$~M\euro{}) : c'est ce rapport, absent de toute comparaison de niveau, qui rend le
transfert pertinent.
\textbf{(b)}~le coût total, prime plus coût du capital résiduel, en fonction de la portée $L$ du
traité en excès de perte ; l'optimum tombe en $L^\star=\mathrm{SCR}=169$~M\euro, là où le
capital résiduel s'annule. (c)~ce que le dimensionnement fait gagner, $-64\,\%$, chiffre
à lire comme une \emph{borne supérieure} : capacité de marché, exclusions et risque de
contrepartie l'entament, mais il faudrait $78\,\%$ d'inefficacité pour renverser la conclusion.}
  \label{fig:detenir}
\end{figure}
% SOURCES-SCRIPTS: 58

\begin{cle}
\textbf{Ce que le raffinement par pilier ne change pas.} Les résultats ci-dessus traitent les
cinq piliers avec une même brique de sévérité et une même détection. Les traiter comme cinq
objets distincts (queue par pilier, accumulation par tiers pour P4, canal de détection
propre) est possible, et c'est l'objet de l'annexe~\ref{chap:adaptations}. Le verdict y est
constant sur les six adaptations testées : chacune relève le niveau, aucune ne
renverse le classement. C'est pourquoi ce matériel est reporté en annexe plutôt que
présenté ici : il renforce la lecture du chapitre sans la modifier.
\end{cle}

% =====================================================================
\section{Les étages d'incertitude et l'auto-excitation : le détail}
\label{sec:ann-etages-hawkes}

% SOURCES-SCRIPTS: 48 71 85

Complément du \S\ref{sec:etage-modele-hors-bande}.

\paragraph{Ce qui justifie le déplacement, et ce n'est pas une taille.} Les deux étages retenus
sont de nature \emph{statistique} : ils mesurent ce que la donnée ne permet pas de trancher à
modèle fixé, et un intervalle est l'objet qui convient pour les porter. Le troisième est
\emph{épistémique} : il mesure l'effet d'un choix de famille, et un choix ne se moyenne pas.

Le point décisif est plus précis que cette opposition générale, et il se voit sur l'axe de la
famille de queue. Cet axe pousse $\xi$ à $0{,}90$, tandis que l'étage de paramètre le fait varier
jusqu'à $0{,}8313$ : c'est le même paramètre, et une bande doit dire \emph{une seule
fois} jusqu'où on le laisse aller. Or $0{,}8313$ est une borne estimée, celle de
l'intervalle bootstrap, quand $0{,}90$ est une valeur posée hors de cet intervalle. Les
placer sur la même échelle met un scénario et une borne de confiance dans le même objet, et le
lecteur ne sait plus si la borne haute est un aléa défavorable ou une autre théorie de la queue.
Le retrait rend la bande homogène, ce qui est un motif plus fort que la lisibilité.

À l'inverse, la borne \emph{basse} de cet axe, la lognormale à $4\,330$~M\euro, est une autre
famille et non un autre $\xi$ : celle-là est un apport propre, et c'est pourquoi le
recouvrement est asymétrique. Le déplacement porte donc sur ce que la bande \emph{annonce},
non sur ce que le mémoire publie.

\paragraph{Où cet étage est publié, désormais.} Il l'est deux fois, et rien n'est caché. Sur
l'axe de la famille de queue, la sixième posture de la table du
\S\ref{sec:predictive-robuste}, le pire cas de modèle à $\xi$ \emph{posé} à $0{,}90$, \emph{est}
l'ambiguïté de famille : c'est le même objet, sous un autre nom, et il porte un chiffre.
Sur l'axe de la dépendance, la bande de $5\,322$ à $9\,806$~M\euro{} reste publiée dans cette
section, avec le repère qui compte, à savoir que la cascade dirigée s'y place entre une
dépendance gaussienne et une dépendance de Gumbel.

\paragraph{Et ce qui ne change pas.} L'ordre d'amorce des piliers reste invariant à la
famille. Le déplacement porte sur la façon de \emph{reporter} le niveau, pas sur la thèse : le
niveau bouge dans la bande, l'ordre résiste, et c'est vrai des trois étages qu'on les cumule
ou non.

\paragraph{Sur l'auto-excitation, une convergence de choix et non une corroboration.} Le parti de
ne pas poser de couche de Hawkes est le plus discutable, l'auto-excitation étant
devenue l'outil par défaut du cyber. Il est ici adossé à un test contre les variantes de noyau de
la littérature. Le préprint de cascade climatique cité au chapitre~\ref{ch:positionnement}
\citep{ccrn2026} arrive au même endroit, et \emph{pour la même raison de résolution temporelle} :
il retient un graphe acyclique à l'échelle de l'événement, au motif que le processus ponctuel
auto-excité devient naturel lorsque les instants exacts et la rétroaction récurrente sont au
c\oe ur du problème, ce qui n'est pas le cas quand la propagation se joue à l'intérieur d'un
événement dont on n'observe pas la chronologie fine. C'est exactement la situation du présent
travail, et la convergence est d'autant plus notable qu'elle est indépendante : autre péril,
autre équipe, et sans connaissance des travaux sur lesquels s'appuie le test conduit ici.

\begin{attention}
\textbf{Ce que cette convergence vaut, et ce qu'elle ne vaut pas.} Il serait commode de la lire
comme une corroboration du rejet ; ce serait la surinterpréter, et la précision importe puisque
c'est le point le plus exposé du mémoire. Ces auteurs ne testent pas l'auto-excitation contre
leur cascade : ils \emph{posent} le graphe acyclique, et leur mention du processus auto-excité
tient en une ligne d'un tableau de positionnement où il figure comme une représentation mature
parmi d'autres, mieux adaptée à un autre régime d'observation. Ce qui converge est donc le
choix de modélisation et son motif, pas une réfutation. Deux équipes qui ne disposent
pas de la chronologie fine d'un événement écartent le même outil pour la même raison, ce qui
déplace ce parti d'une préférence vers une contrainte d'observation, et c'est tout ce que l'on
peut en tirer. Le rejet lui-même ne repose que sur le test du présent travail.
\end{attention}

\paragraph{Ce que ce test établit réellement : une équivalence observationnelle, et c'est plus
fort qu'un rejet.} Le mot \og rejet \fg{} suppose que la donnée départage les deux modèles. Elle
ne les départage pas, et le mesurer valait mieux que de l'espérer.
% SOURCES-SCRIPTS: 85

L'argument se pose d'abord \emph{sans} rien simuler. Si chaque pilier touché est un événement
enregistrable, la fraction d'événements qui sont des \emph{enfants} vaut
$(\mathbb{E}[k]-1)/\mathbb{E}[k]$, où $\mathbb{E}[k]$ est le nombre moyen de piliers touchés par
sinistre : c'est exactement ce qu'un ajustement de Hawkes appelle son ratio de branchement. Avec
$\mathbb{E}[k]=1{,}931$ à l'état non conforme, une cascade sans aucune auto-excitation
prédit $n=0{,}482$, à comparer au $0{,}551$ mesuré sur la chronologie réelle.

Le test le confirme. On engendre soixante années de dates avec le moteur de cascade, dont la
seule structure temporelle est la durée des pas de propagation \emph{à l'intérieur} d'un sinistre,
puis on lui ajuste un Hawkes exponentiel comme on le ferait sur une chronologie réelle. À la
résolution de la donnée disponible, le jour, l'ajustement trouve $n=0{,}480$, soit la prédiction
analytique à trois millièmes. Il retrouve même l'horloge interne : la demi-vie estimée du noyau
vaut $9{,}0$~h pour un lag de propagation vrai de $6{,}7$~h, et il l'appelle \og décroissance de
l'excitation \fg{}. L'estimateur a été validé au préalable sur un Poisson pur, où il rend
$0{,}004$, et sur trois processus de Hawkes simulés de ratios connus, qu'il retrouve.

\begin{cle}
\textbf{Conséquence, à assumer telle quelle.} Un ratio de branchement de l'ordre de $0{,}5$
mesuré sur des dates réelles ne prouve pas l'auto-excitation : une cascade qui n'en
contient aucune le produit. Les deux modèles sont \emph{observationnellement équivalents} à la
résolution disponible. L'énoncé juste n'est donc pas \og le Hawkes est rejeté \fg{} mais \og le
Hawkes n'est pas distinguable de la cascade, et la cascade est retenue par parcimonie
interprétative \fg{} : elle rend compte du même regroupement temporel par un mécanisme qui porte
en outre la structure inter-piliers et se traduit en capital. C'est une conclusion plus faible en
apparence et plus solide en fait, puisqu'elle cesse de s'appuyer sur un départage que la donnée
ne permet pas.

\textbf{Attention au sens de l'effet de résolution.} Dégrader la résolution fait monter
le ratio estimé, de $0{,}473$ à l'heure jusqu'à $0{,}739$ au mois, parce qu'agréger fabrique du
regroupement apparent. C'est l'opération \emph{inverse} de celle conduite plus haut, qui fait
tomber l'endogénéité en reclassant les co-occurrences du même jour en exogène. Les deux
expériences ne se contredisent pas, elles déplacent deux choses différentes, et les confondre
donnerait une conclusion opposée.
\end{cle}

% =====================================================================
\section{Le registre des limites : le détail ligne par ligne}
\label{sec:ann-registre-limites}

% SOURCES-SCRIPTS: 08b 47 57 60 63 65 74 78 89 90 91 102

Complément du \S\ref{sec:pu-coherence}.

La valeur est néanmoins gelée. L'écart représente $2{,}5\,\%$ de la largeur de l'intervalle de
confiance à $90\,\%$ de cette même VaR, $[411{,}5\,;1\,037{,}1]$~M\euro, c'est-à-dire un
quarantième de la précision que le mémoire revendique déjà. Rejouer le pipeline pour le corriger
réinjecterait, sur une grille bootstrap à $200$ tirages, un bruit de simulation du même ordre :
plusieurs centaines de nombres publiés se déplaceraient sans qu'aucun déplacement puisse être
attribué à la correction, et la piste d'audit serait perdue pour un gain immatériel. Le
traitement retenu est celui qu'appelle un défaut de paramètre connu, quantifié, immatériel et
découvert tard, et il vaut mieux que les deux autres postures possibles : on ne recale pas le
modèle, et on ne requalifie pas le couple pour le rendre cohérent sur le papier, ce qui
reviendrait à rendre la réserve moins visible. On inscrit l'écart au registre des limites, avec
son signe et sa taille. Le contrôle qui le mesure est reconduit à chaque exécution, de sorte
qu'il ne peut pas se périmer en silence si le seuil ou le périmètre bougent.

\vspace{2pt}
\noindent \textbf{La troisième ligne est d'une autre taille.} Les deux
premiers défauts pèsent quelques pour cent ; la dérive d'échelle en pèse un quart, soit un ordre
de grandeur de plus, et elle est la seule des trois à avoir été trouvée par un test
\emph{hors échantillon} plutôt que par un contrôle de cohérence interne. Son traitement est
néanmoins le même, et pour le même motif : la corriger supposerait de recalibrer la sévérité sur
une fenêtre glissante, ce qui déplacerait tous les niveaux publiés sans qu'aucun déplacement soit
attribuable à la correction. Deux atténuations la bornent sans l'effacer. La dérive de la queue
vaut le tiers de celle du corps, si bien qu'une lecture rapide de la tendance des pertes
\emph{surestimerait} l'effet d'un facteur trois. Et son effet sur l'\emph{écart} entre états de
conformité est mesuré, non supposé : $-1{,}9\,\%$, soit un déplacement non résolu devant le bruit
de simulation, quand le facteur entre états baisse de $6{,}6\,\%$. La direction est favorable
dans les deux cas, le chiffre publié n'étant pas gonflé par la dérive. Le même backtest valide
par ailleurs la loi de fréquence et la forme de la queue, de sorte que cette ligne est le prix
d'un test dont le solde est favorable (\S\ref{soc:sec:backtest}).

\vspace{2pt}
\noindent \textbf{Et l'argument qui ne doit pas être employé, parce qu'il est faux sous cette
forme.} On serait tenté d'écrire qu'une dérive commune aux deux états se simplifie dans un
rapport, donc que la thèse est à l'abri par construction. La proposition n'est vraie que pour une
variation de la seule \emph{échelle} de sévérité, où elle est d'ailleurs vérifiée,
le rapport des effets sur les deux états valant $0{,}9984$ pour des niveaux qui montent de
$46\,\%$. Elle est \emph{fausse} pour un changement de \emph{forme} de queue, qui déplace le
facteur de $-6{,}1\,\%$, et la dérive mesurée en contient un. La formule juste est donc
l'invariance d'échelle, qui vaut pour toute amplitude, et non l'invariance à la dérive,
qui ne vaut pas.
% SOURCES-SCRIPTS: 90

\vspace{2pt}
\noindent \textbf{Une limite du dispositif de validation, distincte des limites du modèle.} La
même dérive fait rejeter la charge annuelle \emph{agrégée} hors échantillon, ce qui est plus
sévère que sur les marginales puisque l'agrégat cumule la dérive de l'échelle et celle du taux de
dépassement. Aucune limite nouvelle n'entre pour autant à la table, le mécanisme étant le même.
Mais l'exercice révèle une borne de l'instrument qu'il faut déclarer : détecter un taux de
dépassement double du nominal avec une puissance de $80\,\%$ demanderait $1\,811$ années au
niveau de $99{,}5\,\%$. Le quantile qui porte le capital n'est donc backtestable sur aucun
historique de risque opérationnel existant, et ce que le backtest valide est le centre et le
corps de la loi. Le niveau reste adossé à l'ajustement paramétrique et borné par l'intervalle de
$\xi$. Le détail figure au \S\ref{soc:sec:backtest}.
% SOURCES-SCRIPTS: 91

\paragraph{Les deux lignes de périmètre sont d'une autre nature, et vont en sens contraire.} Ce ne
sont ni des défauts d'estimation ni des choix de prudence, mais les deux bornes du périmètre, et
il vaut mieux les nommer ensemble parce qu'elles ne se compensent pas : elles s'additionnent en
incertitude.

La première tient à la donnée. La base de sévérité est tronquée à son seuil de collecte, si bien
que l'attritionnel cyber, les incidents courants à quelques dizaines ou centaines de
milliers d'euros, n'y figure pas du tout. Le besoin de capital publié est donc un capital de
\emph{queue}, et il exclut cette composante non par choix de modèle mais par absence
d'observation. Le sens est clair, il sous-estime, et la taille ne l'est pas : on ne chiffre pas
ce que l'on n'observe pas. Deux remarques en atténuent la portée sans l'annuler. À un niveau de
$99{,}5\,\%$, l'attritionnel pèse peu par construction, puisque le quantile est porté par un
sinistre unique et non par un cumul de petits montants. Et à l'échelle d'une entité, la
fréquence lue vaut $0{,}092$ par an, si bien que $98{,}8\,\%$ des années se passent sans aucun
incident matériel : il n'y a pas, à cette échelle, de régime courant à modéliser. La limite est
réelle au niveau sectoriel, marginale au niveau du quantile d'entité, et elle est ici déclarée
plutôt qu'estimée.

La seconde tient à la transposition. L'élasticité de la sévérité à la taille vaut $0{,}087$, avec
un intervalle qui contient presque zéro : le modèle affirme donc qu'une entité de quelques
milliards d'actifs subit à peu près la sévérité d'une institution mondiale. C'est exactement ce
qui produit la borne inférieure de validité établie au
chapitre~\ref{chap:resultats} et la requalification du niveau d'entité en borne supérieure. Le
correctif est identifié, c'est un plafond de sévérité adossé à l'exposition propre de
l'entité, du même type que celui qui rend le capital PRC calculable lorsque
$\hat\xi>1$. Il n'est pas appliqué, et pour une raison de méthode plutôt que de temps :
l'introduire déplacerait tous les résultats d'entité, ce qui relève de la recalibration et non
de la correction d'erreur. La posture est celle déjà retenue pour $p_u$, une limite écrite et
signée plutôt qu'une reprise silencieuse du pipeline. À noter que cette ligne n'affecte que le
\emph{niveau} d'entité, jamais les écarts entre états de conformité, qui sont des rapports.

\paragraph{Cette limite est désormais chiffrée, et le chiffre est plus parlant que le
coefficient.} Dire que l'élasticité vaut $0{,}087$ énonce la faiblesse en langage de
statisticien. On peut la dire en euros, et c'est le même fait. On plafonne la sévérité d'un
sinistre à une fraction $\kappa$ de l'exposition de l'entité, puis on cherche le $\kappa$ qui
ampute son besoin de moitié.

\begin{center}\footnotesize
\renewcommand{\arraystretch}{1.2}
\begin{tabular}{@{}l r r r r@{}}
\toprule
\textbf{Entité} & \textbf{Exposition} & \textbf{Besoin} & $\boldsymbol{\kappa}$ \textbf{à $-50\,\%$} &
\textbf{soit} \\
\midrule
assureur non-vie A & $2\,319$ & $112{,}1$ & $2{,}284\,\%$ & $53$~M\euro \\
assureur non-vie B & $5\,305$ & $135{,}5$ & $1{,}056\,\%$ & $56$~M\euro \\
assureur vie C & $37\,845$ & $195{,}8$ & $0{,}195\,\%$ & $74$~M\euro \\
assureur vie D & $309\,800$ & $294{,}4$ & \multicolumn{2}{l}{\emph{aucun plafond plausible}} \\
entité notionnelle & $19\,231$ & $169{,}0$ & $0{,}301\,\%$ & $58$~M\euro \\
\bottomrule
\end{tabular}
\end{center}

\vspace{4pt}
\noindent \textbf{Le résultat est dans la dernière colonne.} Rapporté au bilan, le plafond
critique varie d'un facteur $12$ entre entités. \emph{En euros}, il ne va que de $53$ à
$74$~M\euro, un facteur $1{,}4$, sur un panel dont les tailles varient d'un facteur $134$. Le
montant qui plafonne utilement est donc quasiment invariant, ce qui est exactement ce
que l'élasticité de $0{,}087$ affirme, mais dit dans l'unité où un praticien le juge. La
conséquence suit sans statistique : un plafond exprimé en pourcentage du bilan est serré pour une
petite entité et lâche pour une grande.

Et la comparaison qui tranche : pour la plus petite entité du panel, le besoin publié demande
$4{,}83\,\%$ du bilan quand le plafond qui l'amputerait déjà de moitié est à $2{,}28\,\%$. Le
chiffre publié exige donc un plafond deux fois plus lâche que celui qui suffirait à le
couper en deux. Pour la plus grande, le besoin ne demande que $0{,}095\,\%$, loin sous tout
plafond discutable. La borne inférieure de validité se dérive donc d'une hypothèse
économique et non plus d'une propriété de la descente d'échelle : quelle fraction de son bilan
un assureur peut-il perdre sur un seul sinistre informatique. La question se débat en une phrase,
ce qui n'était pas le cas de l'argument qu'elle remplace.
% SOURCES-SCRIPTS: 78

\begin{attention}
\textbf{Un mécanisme découvert en le mesurant, et qui interdit un raccourci.} On attendait que
plafonner chaque sinistre à $\kappa E$ borne le quantile annuel au même niveau, par le principe
de la perte unique dominante. C'est faux dès que le plafond mord : à $\kappa = 0{,}1\,\%$ le
plafond vaut $19$~M\euro{} et le quantile simulé $38$, soit le double. La raison est que
plafonner détruit la queue lourde, donc invalide le principe qui servait à prédire :
la charge annuelle cesse d'être portée par un sinistre et redevient un \emph{cumul}. L'écart va
dans le sens conservateur, le quantile restant au-dessus du plafond, mais un plafond ne fait pas
tomber le capital à $\kappa E$ et il ne faut pas l'annoncer ainsi.
\end{attention}

\paragraph{Et le correctif a désormais une forme fonctionnelle publiée, ce qui change son
statut.} Dire qu'un plafond est \og identifié \fg{} sans en donner la forme laisse le lecteur
juge de la difficulté. Le préprint de cascade climatique cité au chapitre~\ref{ch:positionnement}
\citep{ccrn2026} en propose une, et elle convient au besoin exprimé ici : une perte plafonnée où
la sévérité est bornée par un paramètre d'exposition propre au n\oe ud, assortie d'un effet de
saturation qui fait croître le coût unitaire lorsque plusieurs atteintes surviennent ensemble.
Les deux ingrédients sont exactement ceux qui manquent, le plafond adossé à l'exposition et la
non-linéarité de la simultanéité, cette seconde rejoignant l'hypothèse d'additivité des coûts
discutée au \S\ref{sec:additivite-couts}. Elle n'est pas implémentée, le gel de la calibration
l'interdisant, et aucune valeur de ce préprint n'est reprise, ses paramètres étant
synthétiques. Ce qui change est le statut de la limite : elle cesse d'être une intention pour
devenir une piste référencée, dont la forme est écrite et le coût d'implémentation
connu.

\paragraph{La dernière ligne est la seule dont le sens ne se déclare pas.} Les six précédentes ont
une direction connue, et c'est ce qui permet de les mettre en regard : on sait laquelle sous-estime
et laquelle majore. L'additivité des coûts d'un sinistre multi-piliers n'a pas cette propriété, car
les deux mécanismes qui la mettraient en défaut, la mutualisation de la remédiation et la
saturation de la capacité de réponse, tirent en sens contraire et ne sont pas séparables sur la
donnée disponible. Elle ne peut donc pas entrer dans le bilan des directions, seulement dans celui
des amplitudes, et elle y est la plus grande. Le \S\ref{sec:additivite-couts} la traite pour
elle-même.
